\documentclass[11pt]{report}
% Shared preamble for the Quantum Computing companion lecture notes.
\usepackage[margin=1in]{geometry}
\usepackage{amsmath,amssymb,amsthm,mathtools}
\usepackage{braket}          % \ket, \bra, \braket, \ketbra
\usepackage{enumitem}
\usepackage{graphicx}
\usepackage{xcolor}
\usepackage[colorlinks=true,linkcolor=blue,citecolor=blue,urlcolor=blue]{hyperref}

% ---- Theorem-like environments (shared counter within each chapter) ----
\theoremstyle{definition}
\newtheorem{definition}{Definition}[chapter]
\newtheorem{example}[definition]{Example}
\theoremstyle{plain}
\newtheorem{theorem}[definition]{Theorem}
\newtheorem{lemma}[definition]{Lemma}
\newtheorem{proposition}[definition]{Proposition}
\newtheorem{corollary}[definition]{Corollary}
\theoremstyle{remark}
\newtheorem{remark}[definition]{Remark}

% Exercises use their own per-chapter counter.
\theoremstyle{definition}
\newtheorem{exercise}{Exercise}[chapter]

% Solution environment (unnumbered, italic "Solution." lead-in).
\newenvironment{solution}{\begin{proof}[\normalfont\itshape Solution.]}{\end{proof}}

% ---- Common macros ----
\newcommand{\C}{\mathbb{C}}
\newcommand{\R}{\mathbb{R}}
\newcommand{\Z}{\mathbb{Z}}
\newcommand{\N}{\mathbb{N}}
\newcommand{\F}{\mathbb{F}}
\DeclarePairedDelimiter{\abs}{\lvert}{\rvert}
\DeclarePairedDelimiter{\norm}{\lVert}{\rVert}
\newcommand{\Tr}{\operatorname{Tr}}
\newcommand{\id}{\mathbb{1}}
\newcommand{\CNOT}{\mathrm{CNOT}}
% braket.sty ships \ketbra commented out; define it here so all chapters have it.
\providecommand{\ketbra}[2]{\ket{#1}\!\bra{#2}}
\newcommand{\Had}{H}
\newcommand{\QFT}{\mathrm{QFT}}
\newcommand{\Uf}{U_f}
\newcommand{\defeq}{\vcentcolon=}


\title{\textbf{Quantum Computing}\\[0.3em]\large Companion Lecture Notes}
\author{Companion to the lecture slide decks}
\date{}

\begin{document}
\maketitle

\begin{abstract}
\noindent These notes accompany the six lecture slide decks of the quantum
computing course. They give a self-contained, more detailed treatment of the
same material, with precise mathematical definitions, theorems and lemmas
(with proofs), remarks, worked examples, and problem sets with full solutions.
Notation follows the Dirac convention used in the slides.
\end{abstract}

\tableofcontents

\chapter{Quantum States and Hilbert Space}

% The braket package ships \ketbra commented out; define it if the preamble has not.
\providecommand{\ketbra}[2]{\ket{#1}\!\bra{#2}}

This chapter develops the mathematical stage on which all of quantum computing
takes place. The single postulate we need is deceptively simple: \emph{the state
of a quantum system is a unit vector in a complex inner-product space.} From this
one idea we recover Dirac's bra--ket notation, the geometry of overlaps and
orthogonality, the description of a single qubit and the Bloch sphere, and --- via
the tensor product --- the exponentially large state space of many qubits, where
entanglement first appears. Everything that follows in later chapters (gates,
circuits, and algorithms) is built on top of this vocabulary.

\section{Finite-Dimensional Hilbert Spaces}

\begin{definition}[Complex inner-product space]\label{l1:def:inner}
Let $\mathcal{H}$ be a complex vector space. An \emph{inner product} on
$\mathcal{H}$ is a map $\langle\,\cdot\,,\,\cdot\,\rangle:\mathcal{H}\times\mathcal{H}\to\C$
satisfying, for all $u,v,w\in\mathcal{H}$ and $a,b\in\C$:
\begin{enumerate}[label=(\roman*)]
  \item \emph{Conjugate symmetry:} $\langle u,v\rangle=\overline{\langle v,u\rangle}$.
  \item \emph{Linearity in the second argument:}
        $\langle u,\,a v+b w\rangle=a\langle u,v\rangle+b\langle u,w\rangle$.
        (Combined with (i), this makes the inner product \emph{conjugate-linear}
        in the first argument.)
  \item \emph{Positive-definiteness:} $\langle u,u\rangle\ge 0$, with equality
        if and only if $u=0$.
\end{enumerate}
\end{definition}

We adopt the physicists' convention that the inner product is linear in the
\emph{second} slot and conjugate-linear in the first; this is the convention that
makes Dirac notation come out cleanly. Concretely, on $\C^{N}$ with
$x=(x_0,\dots,x_{N-1})^{T}$ and $y=(y_0,\dots,y_{N-1})^{T}$,
\begin{equation}
  \langle x,y\rangle \;=\; x^{\dagger}y \;=\; \sum_{i=0}^{N-1}\overline{x_i}\,y_i .
\end{equation}

\begin{definition}[Norm and unit vector]\label{l1:def:norm}
The \emph{norm} induced by the inner product is
$\norm{u}=\sqrt{\langle u,u\rangle}$. A vector $u$ is a \emph{unit vector}
(or \emph{normalized}) if $\norm{u}=1$. Two vectors $u,v$ are \emph{orthogonal}
if $\langle u,v\rangle=0$.
\end{definition}

\begin{definition}[Finite-dimensional Hilbert space]\label{l1:def:hilbert}
A \emph{finite-dimensional Hilbert space} is a finite-dimensional complex
inner-product space $\mathcal{H}$ (in finite dimension completeness is automatic,
so no further condition is needed). If $\dim\mathcal{H}=N$, then choosing an
orthonormal basis gives an isomorphism $\mathcal{H}\cong\C^{N}$ that preserves the
inner product. For a system of $n$ qubits we will find $N=2^{n}$.
\end{definition}

The state-space postulate can now be stated precisely.

\begin{remark}[State-space postulate]
The state of an isolated quantum system is described by a \emph{unit} vector in a
Hilbert space $\mathcal{H}$. The space is closed under linear combination: if
$u,v\in\mathcal{H}$ and $a,b\in\C$, then $a u+b v\in\mathcal{H}$. This closure is
the \emph{superposition principle} --- any linear combination of legal states is
again a legal (unnormalized) state.
\end{remark}

\subsection{The inner product controls all geometry}

\begin{proposition}[Elementary consequences of the axioms]\label{l1:prop:elem}
For all $u,v\in\mathcal{H}$ and $a\in\C$:
\begin{enumerate}[label=(\alph*)]
  \item $\langle u,0\rangle=\langle 0,u\rangle=0$;
  \item $\norm{a u}=\abs{a}\,\norm{u}$;
  \item \emph{(Pythagoras)} if $\langle u,v\rangle=0$ then
        $\norm{u+v}^{2}=\norm{u}^{2}+\norm{v}^{2}$.
\end{enumerate}
\end{proposition}

\begin{proof}
(a) By linearity in the second argument, $\langle u,0\rangle=\langle u,0\cdot 0\rangle=0\,\langle u,0\rangle=0$;
the other equality follows from conjugate symmetry.
(b) $\norm{a u}^{2}=\langle a u,a u\rangle=\overline{a}\,a\,\langle u,u\rangle=\abs{a}^{2}\norm{u}^{2}$,
and take square roots.
(c) Expand $\norm{u+v}^{2}=\langle u+v,u+v\rangle=\norm{u}^{2}+\langle u,v\rangle+\langle v,u\rangle+\norm{v}^{2}$;
the two cross terms vanish because $\langle u,v\rangle=0$ forces
$\langle v,u\rangle=\overline{\langle u,v\rangle}=0$.
\end{proof}

The single most useful inequality in the whole subject is the following.

\begin{theorem}[Cauchy--Schwarz inequality]\label{l1:thm:cs}
For all $u,v\in\mathcal{H}$,
\begin{equation}
  \abs{\langle u,v\rangle}\;\le\;\norm{u}\,\norm{v},
\end{equation}
with equality if and only if $u$ and $v$ are linearly dependent.
\end{theorem}

\begin{proof}
If $v=0$ both sides are $0$ and $u,v$ are trivially dependent, so assume $v\neq0$.
For any $\lambda\in\C$, positive-definiteness gives
\[
  0\;\le\;\norm{u-\lambda v}^{2}
   \;=\;\norm{u}^{2}-\overline{\lambda}\langle v,u\rangle-\lambda\langle u,v\rangle+\abs{\lambda}^{2}\norm{v}^{2}.
\]
Choose $\lambda=\dfrac{\langle v,u\rangle}{\norm{v}^{2}}$ (the value that minimizes the
right-hand side). Then $\lambda\langle u,v\rangle=\dfrac{\langle v,u\rangle\langle u,v\rangle}{\norm{v}^{2}}=\dfrac{\abs{\langle u,v\rangle}^{2}}{\norm{v}^{2}}$, and likewise $\overline{\lambda}\langle v,u\rangle=\dfrac{\abs{\langle u,v\rangle}^{2}}{\norm{v}^{2}}$, so
\[
  0\;\le\;\norm{u}^{2}-\frac{\abs{\langle u,v\rangle}^{2}}{\norm{v}^{2}}.
\]
Rearranging yields $\abs{\langle u,v\rangle}^{2}\le\norm{u}^{2}\norm{v}^{2}$, i.e.
the claim. Equality holds iff $\norm{u-\lambda v}=0$, i.e. $u=\lambda v$, which is
exactly linear dependence.
\end{proof}

\begin{corollary}[Triangle inequality]
$\norm{u+v}\le\norm{u}+\norm{v}$ for all $u,v\in\mathcal{H}$.
\end{corollary}

\begin{proof}
$\norm{u+v}^{2}=\norm{u}^{2}+2\,\mathrm{Re}\,\langle u,v\rangle+\norm{v}^{2}
 \le \norm{u}^{2}+2\abs{\langle u,v\rangle}+\norm{v}^{2}
 \le \norm{u}^{2}+2\norm{u}\norm{v}+\norm{v}^{2}=(\norm{u}+\norm{v})^{2}$,
using $\mathrm{Re}\,z\le\abs{z}$ and Theorem~\ref{l1:thm:cs}. Take square roots.
\end{proof}

\subsection{Geometric meaning of the overlap}

For unit vectors the Cauchy--Schwarz inequality reads $\abs{\langle u,v\rangle}\le1$,
so we may write
\begin{equation}\label{l1:eq:geomdecomp}
  \langle u,v\rangle \;=\; \underbrace{\norm{u}\,\norm{v}}_{\text{length}}\;
  \underbrace{\cos\theta}_{\text{alignment}}\;\underbrace{e^{i\varphi}}_{\text{relative phase}},
  \qquad \cos\theta=\frac{\abs{\langle u,v\rangle}}{\norm{u}\,\norm{v}}\in[0,1].
\end{equation}
The modulus measures how much the two directions coincide ($\cos\theta=1$ means
parallel, $\cos\theta=0$ means orthogonal), while the argument $\varphi$ is a
\emph{relative phase} with no analogue for real vectors. Orthogonality,
$\langle u,v\rangle=0$, is exactly the statement that the two states are perfectly
distinguishable.

\begin{example}[An explicit overlap computation]\label{l1:ex:overlap}
Take the unit vectors in $\C^{2}$
\[
  x=\begin{pmatrix}1\\[2pt]0\end{pmatrix},\qquad
  y=\begin{pmatrix}\tfrac{1}{2}e^{i\pi/6}\\[3pt]\tfrac{\sqrt3}{2}\end{pmatrix}.
\]
One checks $\norm{x}^{2}=1$ and $\norm{y}^{2}=\bigl|\tfrac12\bigr|^{2}+\bigl|\tfrac{\sqrt3}{2}\bigr|^{2}=\tfrac14+\tfrac34=1$.
The overlap is
\[
  x^{\dagger}y=\overline{1}\cdot\tfrac12 e^{i\pi/6}+\overline{0}\cdot\tfrac{\sqrt3}{2}
             =\tfrac12 e^{i\pi/6}.
\]
Comparing with \eqref{l1:eq:geomdecomp}: the length factor is $1$, the alignment is
$\cos\theta=\tfrac12$ (so $\theta=60^{\circ}$), and the relative phase is
$e^{i\pi/6}$ (so $\varphi=30^{\circ}$). The angle and the phase are genuinely
different numbers, which is what makes the two factors independent.
\end{example}

\section{Dirac Notation and Quantum States}

\begin{definition}[Ket and bra]\label{l1:def:braket}
A vector of $\mathcal{H}$ is written as a \emph{ket} $\ket{x}$. To each ket there
corresponds a \emph{bra} $\bra{x}$, the element of the dual space
$\mathcal{H}^{*}$ acting by $\bra{x}:\ket{y}\mapsto\langle x,y\rangle$. In
coordinates, if $\ket{x}$ is the column vector with entries $x_i$, then $\bra{x}$
is the row vector $x^{\dagger}=(\overline{x_0},\dots,\overline{x_{N-1}})$, the
conjugate transpose.
\end{definition}

\begin{definition}[Inner and outer products]\label{l1:def:inout}
The \emph{inner product} of $\ket{x}$ and $\ket{y}$ is the scalar
$\braket{x|y}=\langle x,y\rangle\in\C$. The \emph{outer product} $\ketbra{x}{y}$ is
the linear operator acting as
\[
  \bigl(\ketbra{x}{y}\bigr)\ket{z}=\ket{x}\,\braket{y|z}
  =\braket{y|z}\,\ket{x}.
\]
In coordinates it is the matrix with entries $\bigl(\ketbra{x}{y}\bigr)_{ij}=x_i\,\overline{y_j}$.
\end{definition}

Fix an orthonormal basis $\{\ket{i}\}_{i=0}^{N-1}$, meaning $\braket{i|j}=\delta_{ij}$.
Any ket expands as $\ket{x}=\sum_i x_i\ket{i}$, and applying $\bra{j}$ picks out a
single coefficient:
\begin{equation}\label{l1:eq:coeff}
  \braket{j|x}=\sum_i x_i\braket{j|i}=\sum_i x_i\,\delta_{ji}=x_j .
\end{equation}
So the $j$-th amplitude is literally the overlap $\braket{j|x}$.

\begin{theorem}[Completeness / resolution of the identity]\label{l1:thm:completeness}
If $\{\ket{i}\}_{i=0}^{N-1}$ is an orthonormal basis of $\mathcal{H}$, then
\begin{equation}
  \sum_{i=0}^{N-1}\ketbra{i}{i}\;=\;\id .
\end{equation}
Consequently every operator $A$ satisfies $A=\sum_{i,j}A_{ij}\ketbra{i}{j}$ with
$A_{ij}=\braket{i|A|j}$.
\end{theorem}

\begin{proof}
Let $P=\sum_i\ketbra{i}{i}$. For an arbitrary $\ket{x}=\sum_j x_j\ket{j}$,
\[
  P\ket{x}=\sum_i\ket{i}\braket{i|x}=\sum_i x_i\ket{i}=\ket{x},
\]
using \eqref{l1:eq:coeff}. Since $P\ket{x}=\ket{x}$ for every $\ket{x}$, we have
$P=\id$. For the second statement, insert the identity on both sides of $A$:
$A=\id\,A\,\id=\sum_{i,j}\ket{i}\braket{i|A|j}\bra{j}=\sum_{i,j}A_{ij}\ketbra{i}{j}$.
\end{proof}

\begin{remark}[Projectors]
For a unit vector $\ket{u}$ the operator $P_u=\ketbra{u}{u}$ is the orthogonal
\emph{projector} onto the line spanned by $\ket{u}$: it is Hermitian
($P_u^{\dagger}=P_u$) and idempotent, since
$P_u^{2}=\ket{u}\braket{u|u}\bra{u}=\ket{u}\bra{u}=P_u$ because $\braket{u|u}=1$.
The completeness relation of Theorem~\ref{l1:thm:completeness} decomposes $\id$ into a
sum of orthogonal rank-one projectors, one per basis vector.
\end{remark}

\subsection{The single qubit}

\begin{definition}[Qubit and computational basis]\label{l1:def:qubit}
A \emph{qubit} is a quantum system whose state space is $\mathcal{H}=\C^{2}$. The
distinguished orthonormal basis
\[
  \ket{0}=\begin{pmatrix}1\\0\end{pmatrix},\qquad
  \ket{1}=\begin{pmatrix}0\\1\end{pmatrix}
\]
is called the \emph{computational basis}. A general pure state is
\begin{equation}\label{l1:eq:qubit}
  \ket{\psi}=a\ket{0}+b\ket{1},\qquad a,b\in\C,\qquad \abs{a}^{2}+\abs{b}^{2}=1 .
\end{equation}
The coefficients $a,b$ are called \emph{amplitudes}; $\abs{a}^{2}$ and $\abs{b}^{2}$
are the probabilities of obtaining $0$ or $1$ in a computational-basis measurement.
\end{definition}

\begin{remark}[Global phase is unobservable]\label{l1:rem:global}
Measurement probabilities depend on the amplitudes only through moduli of the form
$\abs{\braket{i|\psi}}^{2}$. Replacing $\ket{\psi}$ by $e^{i\gamma}\ket{\psi}$ for a
real \emph{global phase} $\gamma$ leaves every such probability unchanged, and more
generally leaves $\braket{\phi|A|\psi}$-type quantities invariant when the same
phase appears on both sides. Hence $e^{i\gamma}\ket{\psi}$ and $\ket{\psi}$ describe
the same physical state: genuine states are points of the \emph{projective space}
$\mathbb{P}(\mathcal{H})$, i.e. rays rather than vectors. \emph{Relative} phases
between the terms of a superposition, by contrast, are physical and drive
interference.
\end{remark}

\begin{proposition}[Bloch parametrization]\label{l1:prop:bloch}
Every single-qubit pure state can be written, up to an unobservable global phase, as
\begin{equation}\label{l1:eq:blochform}
  \ket{\psi}=\cos\tfrac{\theta}{2}\,\ket{0}+e^{i\varphi}\sin\tfrac{\theta}{2}\,\ket{1},
  \qquad \theta\in[0,\pi],\quad \varphi\in[0,2\pi).
\end{equation}
The pair $(\theta,\varphi)$ is the pair of spherical angles of a point on the unit
sphere $S^{2}\subset\R^{3}$, the \emph{Bloch sphere}.
\end{proposition}

\begin{proof}
Start from \eqref{l1:eq:qubit} with $\abs{a}^{2}+\abs{b}^{2}=1$. Write the amplitudes in
polar form $a=r_a e^{i\alpha}$ and $b=r_b e^{i\beta}$ with $r_a,r_b\ge0$. Since
$r_a^{2}+r_b^{2}=1$ we may set $r_a=\cos\tfrac{\theta}{2}$ and
$r_b=\sin\tfrac{\theta}{2}$ for a unique $\theta\in[0,\pi]$ (both cosine and sine of
$\theta/2$ are then nonnegative). Thus
\[
  \ket{\psi}=e^{i\alpha}\cos\tfrac{\theta}{2}\ket{0}+e^{i\beta}\sin\tfrac{\theta}{2}\ket{1}
           =e^{i\alpha}\Bigl(\cos\tfrac{\theta}{2}\ket{0}+e^{i(\beta-\alpha)}\sin\tfrac{\theta}{2}\ket{1}\Bigr).
\]
The overall factor $e^{i\alpha}$ is a global phase, which by
Remark~\ref{l1:rem:global} we may discard. Setting $\varphi=\beta-\alpha\ \mathrm{mod}\ 2\pi$
gives \eqref{l1:eq:blochform}. When $\theta\in\{0,\pi\}$ the state is $\ket{0}$ or
$\ket{1}$ and $\varphi$ is irrelevant (the poles of the sphere), consistent with the
identification of states with points of $S^{2}$.
\end{proof}

\begin{remark}[Reading the Bloch sphere]
In this picture $\ket{0}$ and $\ket{1}$ sit at the north and south poles
($\theta=0,\pi$), and \emph{orthogonal} qubit states are \emph{antipodal} points,
not perpendicular ones --- a state and its orthogonal complement are separated by
angle $\theta=\pi$ on the sphere even though they are separated by $90^{\circ}$ in
$\C^{2}$. The equatorial states $\tfrac{1}{\sqrt2}(\ket{0}+e^{i\varphi}\ket{1})$ have
$\theta=\pi/2$ and differ only in their relative phase $\varphi$.
\end{remark}

\begin{example}[Amplitudes and normalization]
The state $\ket{+}=\tfrac{1}{\sqrt2}(\ket{0}+\ket{1})$ has
$\abs{a}^{2}=\abs{b}^{2}=\tfrac12$, so a computational-basis measurement yields $0$ or
$1$ with equal probability $\tfrac12$; its Bloch angles are
$\theta=\pi/2,\varphi=0$. The state $\ket{-}=\tfrac{1}{\sqrt2}(\ket{0}-\ket{1})$ has
the \emph{same} measurement statistics in the computational basis but differs by the
relative phase $\varphi=\pi$; indeed $\braket{+|-}=\tfrac12(1-1)=0$, so $\ket{+}$ and
$\ket{-}$ are orthogonal and perfectly distinguishable by a suitable measurement.
\end{example}

\section{Multi-Qubit Systems and Entanglement}

\begin{definition}[Tensor-product state space]\label{l1:def:tensor}
If systems $A$ and $B$ have state spaces $\mathcal{H}_A$ and $\mathcal{H}_B$, the
composite system has state space the \emph{tensor product}
$\mathcal{H}_A\otimes\mathcal{H}_B$. If $\{\ket{i}_A\}$ and $\{\ket{j}_B\}$ are
orthonormal bases, then $\{\ket{i}_A\otimes\ket{j}_B\}$ is an orthonormal basis of the
product, with inner product
\[
  \bigl(\bra{i}\otimes\bra{j}\bigr)\bigl(\ket{k}\otimes\ket{l}\bigr)
    =\braket{i|k}\,\braket{j|l}.
\]
Hence $\dim(\mathcal{H}_A\otimes\mathcal{H}_B)=(\dim\mathcal{H}_A)(\dim\mathcal{H}_B)$.
For $n$ qubits,
\[
  \mathcal{H}=(\C^{2})^{\otimes n}\cong\C^{2^{n}},\qquad \dim\mathcal{H}=2^{n}.
\]
\end{definition}

We abbreviate $\ket{i}\otimes\ket{j}$ as $\ket{i}\ket{j}$ or $\ket{ij}$. The
computational basis of $n$ qubits is $\{\ket{x}:x\in\{0,1\}^{n}\}$, indexed by the
$2^{n}$ bit strings. The exponential growth of the dimension --- $50$ qubits already
give a space of dimension $2^{50}>10^{15}$ --- is one root of quantum computing's
potential power: a classical description of a general state requires exponentially
many amplitudes.

\begin{definition}[Product and entangled states]\label{l1:def:product}
A state $\ket{\psi}\in\mathcal{H}_A\otimes\mathcal{H}_B$ is a \emph{product state}
(or \emph{separable}) if it can be written as $\ket{\psi}=\ket{\alpha}_A\otimes\ket{\beta}_B$
for some $\ket{\alpha}_A\in\mathcal{H}_A$, $\ket{\beta}_B\in\mathcal{H}_B$. Otherwise
$\ket{\psi}$ is \emph{entangled}.
\end{definition}

\begin{example}[A generic two-qubit state]
Two qubits live in $\C^{4}$ with basis $\ket{00},\ket{01},\ket{10},\ket{11}$. A
general state is
\[
  \ket{\psi}=\sum_{i,j\in\{0,1\}}x_{ij}\ket{ij},\qquad
  \sum_{i,j}\abs{x_{ij}}^{2}=1 .
\]
For instance $\tfrac12(\ket{00}+\ket{01}+\ket{10}+\ket{11})$ is normalized
($4\cdot\tfrac14=1$) and, as the reader can check, factors as
$\tfrac{1}{\sqrt2}(\ket{0}+\ket{1})\otimes\tfrac{1}{\sqrt2}(\ket{0}+\ket{1})$, so it
is a product state.
\end{example}

\begin{definition}[Bell state]\label{l1:def:bell}
The \emph{Bell state} (or EPR pair) is the two-qubit state
\begin{equation}
  \ket{B^{+}}=\tfrac{1}{\sqrt2}\bigl(\ket{00}+\ket{11}\bigr).
\end{equation}
\end{definition}

\begin{theorem}[The Bell state is entangled]\label{l1:thm:bell}
There is no pair of single-qubit states $\ket{\alpha},\ket{\beta}\in\C^{2}$ with
$\ket{B^{+}}=\ket{\alpha}\otimes\ket{\beta}$.
\end{theorem}

\begin{proof}
Suppose for contradiction that such a factorization exists. Write
$\ket{\alpha}=a_0\ket{0}+a_1\ket{1}$ and $\ket{\beta}=b_0\ket{0}+b_1\ket{1}$. Then
\[
  \ket{\alpha}\otimes\ket{\beta}
    =a_0 b_0\ket{00}+a_0 b_1\ket{01}+a_1 b_0\ket{10}+a_1 b_1\ket{11}.
\]
Matching coefficients with $\ket{B^{+}}=\tfrac{1}{\sqrt2}\ket{00}+0\cdot\ket{01}+0\cdot\ket{10}+\tfrac{1}{\sqrt2}\ket{11}$ requires
\[
  a_0 b_0=\tfrac{1}{\sqrt2},\qquad a_0 b_1=0,\qquad a_1 b_0=0,\qquad a_1 b_1=\tfrac{1}{\sqrt2}.
\]
From $a_0 b_0=\tfrac{1}{\sqrt2}\neq0$ we get $a_0\neq0$ and $b_0\neq0$. Then
$a_0 b_1=0$ forces $b_1=0$, while $a_1 b_0=0$ forces $a_1=0$; hence
$a_1 b_1=0\neq\tfrac{1}{\sqrt2}$, a contradiction. Therefore no factorization exists
and $\ket{B^{+}}$ is entangled.
\end{proof}

\begin{remark}[Physical meaning of entanglement]
The correlations of $\ket{B^{+}}$ have no separable description. Measuring the first
qubit in the computational basis yields $0$ or $1$ with probability $\tfrac12$ each;
whatever the outcome, the post-measurement state of the second qubit is forced to the
\emph{same} value ($0\mapsto\ket{0}$, $1\mapsto\ket{1}$). The two outcomes are
perfectly correlated even if the qubits are spatially far apart. This correlation,
inexpressible as any classical mixture of independent single-qubit states, is the
resource behind teleportation, superdense coding, and much of quantum advantage.
\end{remark}

\section{Problems}

\begin{exercise}\label{l1:ex:norm}
Let $\ket{\psi}=2i\ket{0}-\ket{1}\in\C^{2}$ (unnormalized). Find $\norm{\psi}$ and
the normalized state $\ket{\hat\psi}$, and give the probabilities of measuring $0$
and $1$ in the computational basis.
\end{exercise}

\begin{exercise}\label{l1:ex:orth}
Show that $\ket{+}=\tfrac{1}{\sqrt2}(\ket{0}+\ket{1})$ and
$\ket{-}=\tfrac{1}{\sqrt2}(\ket{0}-\ket{1})$ form an orthonormal basis of $\C^{2}$,
and express $\ket{0}$ and $\ket{1}$ in this basis.
\end{exercise}

\begin{exercise}\label{l1:ex:cseq}
Using the equality case of the Cauchy--Schwarz inequality
(Theorem~\ref{l1:thm:cs}), prove that two \emph{unit} vectors $\ket{u},\ket{v}$ satisfy
$\abs{\braket{u|v}}=1$ if and only if $\ket{v}=e^{i\gamma}\ket{u}$ for some real
$\gamma$ --- i.e. they are the same physical state.
\end{exercise}

\begin{exercise}\label{l1:ex:proj}
Let $\ket{u}$ be a unit vector and $P=\ketbra{u}{u}$. Prove that $P^{2}=P$ and
$P^{\dagger}=P$, and compute $\Tr(P)$.
\end{exercise}

\begin{exercise}\label{l1:ex:bloch}
Find the Bloch angles $(\theta,\varphi)$ of the state
$\ket{\psi}=\tfrac{1}{2}\ket{0}+\tfrac{\sqrt3}{2}e^{i\pi/3}\ket{1}$.
\end{exercise}

\begin{exercise}\label{l1:ex:sep}
Decide whether $\ket{\psi}=\tfrac{1}{\sqrt2}(\ket{00}+\ket{01})$ is a product state.
If so, exhibit an explicit factorization; if not, prove it is entangled.
\end{exercise}

\begin{exercise}\label{l1:ex:bellmeas}
For the Bell state $\ket{B^{+}}=\tfrac{1}{\sqrt2}(\ket{00}+\ket{11})$, compute the
probability of obtaining outcome $0$ on the \emph{first} qubit, and give the
resulting (normalized) state of the second qubit.
\end{exercise}

\section*{Solutions}

\begin{solution}[Exercise~\ref{l1:ex:norm}]
$\norm{\psi}^{2}=\abs{2i}^{2}+\abs{-1}^{2}=4+1=5$, so $\norm{\psi}=\sqrt5$. The
normalized state is $\ket{\hat\psi}=\tfrac{1}{\sqrt5}(2i\ket{0}-\ket{1})$. The
probabilities are $\abs{2i/\sqrt5}^{2}=4/5$ for outcome $0$ and
$\abs{-1/\sqrt5}^{2}=1/5$ for outcome $1$; they sum to $1$ as required.
\end{solution}

\begin{solution}[Exercise~\ref{l1:ex:orth}]
$\braket{+|+}=\tfrac12(1+1)=1$ and $\braket{-|-}=\tfrac12(1+1)=1$, so both are unit
vectors. $\braket{+|-}=\tfrac12(\braket{0|0}-\braket{0|1}+\braket{1|0}-\braket{1|1})=\tfrac12(1-0+0-1)=0$,
so they are orthogonal; two orthonormal vectors in a $2$-dimensional space form a
basis. Adding and subtracting, $\ket{0}=\tfrac{1}{\sqrt2}(\ket{+}+\ket{-})$ and
$\ket{1}=\tfrac{1}{\sqrt2}(\ket{+}-\ket{-})$.
\end{solution}

\begin{solution}[Exercise~\ref{l1:ex:cseq}]
For unit vectors, $\norm{u}=\norm{v}=1$, so Cauchy--Schwarz gives
$\abs{\braket{u|v}}\le1$ with equality iff $\ket{u},\ket{v}$ are linearly dependent,
i.e. $\ket{v}=\lambda\ket{u}$ for some $\lambda\in\C$. Taking norms,
$1=\norm{v}=\abs{\lambda}\norm{u}=\abs{\lambda}$, so $\lambda=e^{i\gamma}$ for some
real $\gamma$. Conversely, if $\ket{v}=e^{i\gamma}\ket{u}$ then
$\braket{u|v}=e^{i\gamma}$ and $\abs{\braket{u|v}}=1$. By Remark~\ref{l1:rem:global} the
two are the same physical state.
\end{solution}

\begin{solution}[Exercise~\ref{l1:ex:proj}]
$P^{2}=\ket{u}\braket{u|u}\bra{u}=\ket{u}\bra{u}=P$ since $\braket{u|u}=1$. Taking the
adjoint, $P^{\dagger}=(\ketbra{u}{u})^{\dagger}=\ketbra{u}{u}=P$. For the trace,
using a basis extension or the identity $\Tr(\ketbra{u}{v})=\braket{v|u}$, we get
$\Tr(P)=\braket{u|u}=1$ (a rank-one projector has trace equal to its rank).
\end{solution}

\begin{solution}[Exercise~\ref{l1:ex:bloch}]
Compare with \eqref{l1:eq:blochform}. The $\ket{0}$ amplitude is real and nonnegative,
$\cos\tfrac{\theta}{2}=\tfrac12$, giving $\theta/2=\pi/3$, so $\theta=2\pi/3$. Then
$\sin\tfrac{\theta}{2}=\sin\tfrac{\pi}{3}=\tfrac{\sqrt3}{2}$, matching the modulus of
the $\ket{1}$ amplitude, and its phase is $\varphi=\pi/3$. Hence
$(\theta,\varphi)=(2\pi/3,\ \pi/3)$.
\end{solution}

\begin{solution}[Exercise~\ref{l1:ex:sep}]
Factor out the first qubit:
$\ket{\psi}=\tfrac{1}{\sqrt2}\ket{0}\otimes(\ket{0}+\ket{1})
=\ket{0}\otimes\tfrac{1}{\sqrt2}(\ket{0}+\ket{1})=\ket{0}\otimes\ket{+}$.
This is an explicit product of single-qubit states, so $\ket{\psi}$ is separable
(a product state), not entangled.
\end{solution}

\begin{solution}[Exercise~\ref{l1:ex:bellmeas}]
Write $\ket{B^{+}}=\tfrac{1}{\sqrt2}\ket{0}\ket{0}+\tfrac{1}{\sqrt2}\ket{1}\ket{1}$.
The projector for outcome $0$ on the first qubit is $\ketbra{0}{0}\otimes\id$. The
outcome probability is the squared norm of the (unnormalized) post-measurement vector
$(\ketbra{0}{0}\otimes\id)\ket{B^{+}}=\tfrac{1}{\sqrt2}\ket{0}\ket{0}$, namely
$\bigl\|\tfrac{1}{\sqrt2}\ket{00}\bigr\|^{2}=\tfrac12$. After normalization the state
is $\ket{0}\otimes\ket{0}$, so the second qubit is left in $\ket{0}$ --- perfectly
correlated with the first outcome, as expected for a Bell pair.
\end{solution}

\chapter{Quantum Gates, Circuits, and the No-Cloning Theorem}

This chapter develops the operational layer of quantum computing. We treat gates
as unitary operators, catalogue the standard single- and multi-qubit gates and
the algebraic identities relating them, learn to read circuit diagrams, discuss
universal gate sets and compilation to hardware, and finally state and prove the
no-cloning theorem together with its consequences.

\section{Unitary Operators as Gates}

The postulates of quantum mechanics say that the state of an isolated system
evolves in time by a linear, inner-product-preserving map. On a finite
dimensional Hilbert space such a map is exactly a unitary operator, and in
quantum computing we call it a gate.

\begin{definition}[Unitary operator / quantum gate]
Let $\mathcal{H}\cong\C^{d}$ be a finite-dimensional Hilbert space. A linear
operator $U:\mathcal{H}\to\mathcal{H}$ is \emph{unitary} if
\[
  U^{\dagger}U = U U^{\dagger} = \id,
\]
where $U^{\dagger}$ is the conjugate transpose (adjoint) of $U$. A
\emph{quantum gate} on $n$ qubits is a unitary operator on
$\mathcal{H}=(\C^{2})^{\otimes n}\cong\C^{2^{n}}$.
\end{definition}

Equivalently, $U$ is unitary iff its columns (or rows) form an orthonormal basis
of $\C^{d}$, iff $U^{-1}=U^{\dagger}$. In particular every gate is invertible, so
quantum computation is reversible.

\begin{theorem}[Gates preserve inner products and norm]
\label{l2:thm:unitary-preserves}
Let $U$ be a linear operator on $\mathcal{H}$. The following are equivalent.
\begin{enumerate}[label=\textup{(\roman*)}]
  \item $U$ is unitary, i.e. $U^{\dagger}U=\id$.
  \item $U$ preserves inner products: $\braket{U\psi|U\phi}=\braket{\psi|\phi}$
        for all $\ket{\psi},\ket{\phi}\in\mathcal{H}$.
  \item $U$ preserves norm: $\norm{U\psi}=\norm{\psi}$ for all
        $\ket{\psi}\in\mathcal{H}$.
\end{enumerate}
\end{theorem}

\begin{proof}
$(\mathrm{i})\Rightarrow(\mathrm{ii})$. Using
$\braket{U\psi|U\phi}=\bra{\psi}U^{\dagger}U\ket{\phi}$ and $U^{\dagger}U=\id$,
we get $\braket{U\psi|U\phi}=\bra{\psi}\id\ket{\phi}=\braket{\psi|\phi}$.

$(\mathrm{ii})\Rightarrow(\mathrm{iii})$. Take $\ket{\phi}=\ket{\psi}$: then
$\norm{U\psi}^{2}=\braket{U\psi|U\psi}=\braket{\psi|\psi}=\norm{\psi}^{2}$, and
taking nonnegative square roots gives $\norm{U\psi}=\norm{\psi}$.

$(\mathrm{iii})\Rightarrow(\mathrm{i})$. Norm preservation gives
$\bra{\psi}U^{\dagger}U\ket{\psi}=\bra{\psi}\id\ket{\psi}$ for all $\ket{\psi}$,
so the self-adjoint operator $M\defeq U^{\dagger}U-\id$ satisfies
$\bra{\psi}M\ket{\psi}=0$ for every $\ket{\psi}$. For a self-adjoint operator the
quadratic form determines the operator (via the polarization identity, or by
choosing $\ket{\psi}$ to be its eigenvectors), hence $M=0$, i.e.
$U^{\dagger}U=\id$. On a finite-dimensional space a left inverse is a two-sided
inverse, so also $UU^{\dagger}=\id$ and $U$ is unitary.
\end{proof}

Physically, Theorem~\ref{l2:thm:unitary-preserves} says gates map normalized states
to normalized states and never change the ``angle'' between two states; this
inner-product rigidity is exactly what will forbid cloning in
Section~\ref{l2:sec:nocloning}.

\section{Single-Qubit Gates}

\begin{definition}[Pauli, Hadamard, and phase gates]
The \emph{Pauli gates} and identity on $\C^{2}$ are
\[
  I=\begin{pmatrix}1&0\\0&1\end{pmatrix},\quad
  X=\begin{pmatrix}0&1\\1&0\end{pmatrix},\quad
  Y=\begin{pmatrix}0&-i\\i&0\end{pmatrix},\quad
  Z=\begin{pmatrix}1&0\\0&-1\end{pmatrix}.
\]
The \emph{Hadamard gate} and the \emph{phase gates} $S$ and $T$ are
\[
  H=\frac{1}{\sqrt{2}}\begin{pmatrix}1&1\\1&-1\end{pmatrix},\qquad
  S=\begin{pmatrix}1&0\\0&i\end{pmatrix},\qquad
  T=\begin{pmatrix}1&0\\0&e^{i\pi/4}\end{pmatrix}.
\]
\end{definition}

Each of these is unitary: one checks directly that $X^{\dagger}X=Y^{\dagger}Y
=Z^{\dagger}Z=I$, that $H^{\dagger}H=I$, and that $S,T$ are diagonal with
unit-modulus entries. In the computational basis,
\[
  X\ket{0}=\ket{1},\quad X\ket{1}=\ket{0},\qquad
  Z\ket{0}=\ket{0},\quad Z\ket{1}=-\ket{1},
\]
so $X$ is a \emph{bit flip} (the quantum \textsc{not} gate) and $Z$ is a
\emph{phase flip}: it fixes $\ket{0}$ and multiplies $\ket{1}$ by $-1$. The
Hadamard gate creates superpositions,
\[
  H\ket{0}=\ket{+}=\tfrac{1}{\sqrt{2}}(\ket{0}+\ket{1}),\qquad
  H\ket{1}=\ket{-}=\tfrac{1}{\sqrt{2}}(\ket{0}-\ket{1}),
\]
and is self-inverse, $H^{2}=I$, since
$\tfrac12\left(\begin{smallmatrix}1&1\\1&-1\end{smallmatrix}\right)
\left(\begin{smallmatrix}1&1\\1&-1\end{smallmatrix}\right)
=\left(\begin{smallmatrix}1&0\\0&1\end{smallmatrix}\right)$. The phase gates
leave $\ket{0}$ alone and rotate the phase of $\ket{1}$: $S$ multiplies it by
$i$ and $T$ by $e^{i\pi/4}$.

The four operators $\{I,X,Y,Z\}$ form a basis of the (complex) vector space of
$2\times2$ matrices, so every single-qubit gate can be written as a complex
linear combination of Pauli matrices. Geometrically, a single-qubit pure state
corresponds to a point on the \emph{Bloch sphere}; the Pauli gates act as
rotations by $\pi$ about the $x$, $y$, $z$ axes (up to a global phase), $X$ being
a bit flip about the $x$-axis and $Z$ a phase flip about the $z$-axis.

\subsection{Algebraic identities}

The Pauli, Hadamard, and phase gates satisfy several identities used constantly
in practice. We prove the two conjugation identities that exchange bit and phase
flips.

\begin{lemma}[Hadamard conjugation]
\label{l2:lem:conj}
$H X H = Z$ and $H Z H = X$.
\end{lemma}

\begin{proof}
Recall $H=H^{\dagger}$ and $H^{2}=I$. Compute $HXH$ by matrix multiplication.
First
\[
  XH=\begin{pmatrix}0&1\\1&0\end{pmatrix}
     \frac{1}{\sqrt2}\begin{pmatrix}1&1\\1&-1\end{pmatrix}
    =\frac{1}{\sqrt2}\begin{pmatrix}1&-1\\1&1\end{pmatrix},
\]
and then
\[
  HXH=\frac{1}{\sqrt2}\begin{pmatrix}1&1\\1&-1\end{pmatrix}\cdot
      \frac{1}{\sqrt2}\begin{pmatrix}1&-1\\1&1\end{pmatrix}
     =\frac{1}{2}\begin{pmatrix}2&0\\0&-2\end{pmatrix}
     =\begin{pmatrix}1&0\\0&-1\end{pmatrix}=Z.
\]
For the second identity, conjugate the first by $H$ on both sides and use
$H^{2}=I$: from $HXH=Z$ we get $H(HXH)H=HZH$, i.e. $(H^{2})X(H^{2})=HZH$, so
$X=HZH$. Alternatively compute directly:
\[
  ZH=\frac{1}{\sqrt2}\begin{pmatrix}1&1\\-1&1\end{pmatrix},\qquad
  HZH=\frac{1}{2}\begin{pmatrix}1&1\\1&-1\end{pmatrix}
      \begin{pmatrix}1&1\\-1&1\end{pmatrix}
     =\begin{pmatrix}0&1\\1&0\end{pmatrix}=X. \qedhere
\]
\end{proof}

Lemma~\ref{l2:lem:conj} is the precise statement that $H$ is a ``bridge'' between
the bit-flip and phase-flip pictures: a phase flip in the computational basis
looks like a bit flip in the Hadamard ($\ket{+},\ket{-}$) basis, and vice versa.
Two further identities are recorded for reference.

\begin{proposition}
\label{l2:prop:yxz}
$Y=iXZ$ and $S=T^{2}$.
\end{proposition}

\begin{proof}
Direct computation:
\[
  XZ=\begin{pmatrix}0&1\\1&0\end{pmatrix}
     \begin{pmatrix}1&0\\0&-1\end{pmatrix}
    =\begin{pmatrix}0&-1\\1&0\end{pmatrix},\qquad
  iXZ=\begin{pmatrix}0&-i\\i&0\end{pmatrix}=Y.
\]
For the phase gates, $T$ is diagonal so
$T^{2}=\operatorname{diag}(1,e^{i\pi/4})^{2}=\operatorname{diag}(1,e^{i\pi/2})
=\operatorname{diag}(1,i)=S$.
\end{proof}

\begin{remark}[$T$ is non-Clifford]
The \emph{Clifford group} is generated by $\{H,S,\CNOT\}$ (equivalently, the
unitaries that map the Pauli group to itself under conjugation). The identities
above show $X,Y,Z$ and $S=T^{2}$ lie in this group. The $T$ gate does
\emph{not}: conjugating a Pauli by $T$ produces a superposition of Paulis rather
than another Pauli, so $T$ is a non-Clifford gate. This matters because the
Gottesman--Knill theorem says Clifford circuits are efficiently classically
simulable; the $T$ gate is precisely the ingredient that lifts a gate set to
universality for quantum computation.
\end{remark}

\section{Multi-Qubit Gates}

To act on more than one qubit we take tensor products. A gate $A$ on qubit~$0$
and $B$ on qubit~$1$ acts on the joint state as $A\otimes B$, with
$(A\otimes B)(\ket{u}\otimes\ket{v})=A\ket{u}\otimes B\ket{v}$. The most
important genuinely two-qubit gate is the controlled-\textsc{not}.

\begin{definition}[$\CNOT$ and Toffoli]
The \emph{controlled-\textsc{not}} gate $\CNOT$ acts on two qubits (control
$a$, target $b$) by
\[
  \CNOT\ket{a,b}=\ket{a,\,a\oplus b},\qquad a,b\in\{0,1\},
\]
where $\oplus$ is addition modulo $2$. In the ordered computational basis
$\ket{00},\ket{01},\ket{10},\ket{11}$ its matrix is the $4\times4$ permutation
unitary
\[
  \CNOT=\begin{pmatrix}
    1&0&0&0\\ 0&1&0&0\\ 0&0&0&1\\ 0&0&1&0
  \end{pmatrix}.
\]
The \emph{Toffoli} (or \textsc{ccnot}) gate acts on three qubits by
\[
  \mathrm{CCNOT}\ket{a,b,c}=\ket{a,\,b,\,(a\wedge b)\oplus c},
\]
flipping the target $c$ exactly when both controls $a,b$ equal $1$.
\end{definition}

Since $\CNOT$ permutes the four basis vectors it is a permutation matrix, hence
unitary; the same holds for the Toffoli gate, which permutes the eight
three-qubit basis vectors by swapping $\ket{110}\leftrightarrow\ket{111}$. Both
are their own inverse. The Toffoli gate is universal for classical reversible
computation: with ancillas it can implement \textsc{and}, \textsc{not}, and
fan-out, so any Boolean circuit can be made reversible.

\begin{remark}[Why XOR keeps oracles reversible]
A classical oracle for $f:\{0,1\}^{n}\to\{0,1\}$ is realized quantumly as
$\Uf\ket{x}\ket{b}=\ket{x}\ket{b\oplus f(x)}$. Writing the answer with $\oplus$
rather than overwriting $b$ with $f(x)$ is what makes $\Uf$ a permutation of
basis states, hence unitary and reversible: applying $\Uf$ twice returns
$\ket{b\oplus f(x)\oplus f(x)}=\ket{b}$, so $\Uf^{2}=\id$. A plain overwrite
$b\mapsto f(x)$ is generally not injective and cannot be unitary. The $\CNOT$
rule $b\mapsto a\oplus b$ is the single-bit instance of this construction.
\end{remark}

\section{Reading Quantum Circuits}

A \emph{quantum circuit} is a diagram in which each horizontal wire is a qubit
and gates are drawn on the wires. Three conventions govern how to read one.

\begin{enumerate}[label=(\arabic*)]
  \item \textbf{Time flows left to right.} The input state sits at the far left;
        a gate drawn further left is applied earlier.
  \item \textbf{Controls.} On a controlled gate, a filled dot $\bullet$ marks a
        control that activates when the qubit is $\ket{1}$; an open dot
        $\circ$ marks a $0$-control, which activates when the qubit is $\ket{0}$.
  \item \textbf{Matrix order is reversed.} If the circuit applies $U_{1}$ first,
        then $U_{2}$, the overall unitary is the \emph{reverse-order} product
        \[
          U = U_{2}\,U_{1},
        \]
        because operators act on the ket to their right. For a length-$m$
        sequence, $U=U_{m}U_{m-1}\cdots U_{1}$.
\end{enumerate}

A controlled-$U$ gate with control $a$ can be written compactly as
$c\text{-}U\,\ket{a}\ket{b}=\ket{a}\,U^{a}\ket{b}$, so $U$ is applied to the
target raised to the power $a\in\{0,1\}$ (nothing happens when $a=0$). The
\textsc{swap} gate exchanges two qubits, $\mathrm{SWAP}\ket{a,b}=\ket{b,a}$; it
can be built from three $\CNOT$s.

\begin{definition}[Circuit depth]
Group the gates of a circuit into \emph{layers}, where a layer is a maximal set
of gates acting on disjoint qubits that can be executed simultaneously. The
\emph{depth} of the circuit is the number of layers, equivalently the largest
number of gates on any single input-to-output path through the diagram.
\end{definition}

Gates acting on disjoint sets of qubits commute as tensor factors and run in
parallel, so depth (not total gate count) measures wall-clock runtime, and
keeping depth low is central to running on noisy hardware.

\begin{example}[A small circuit's unitary as a reverse-order product]
Consider one qubit with, from left to right, an $H$ gate, then an $S$ gate, then
another $H$ gate. The overall unitary is the reverse-order product
\[
  U = H\,S\,H
    = \frac{1}{2}\begin{pmatrix}1&1\\1&-1\end{pmatrix}
      \begin{pmatrix}1&0\\0&i\end{pmatrix}
      \begin{pmatrix}1&1\\1&-1\end{pmatrix}
    = \frac{1}{2}\begin{pmatrix}1+i&1-i\\1-i&1+i\end{pmatrix}.
\]
One checks $U^{\dagger}U=I$, confirming the composite is again a gate. Reading
the diagram in the wrong (left-to-right) matrix order would instead compute
$HSH$ with the factors reversed, which here gives the same matrix only because
the outer gates coincide; in general the two orders differ.
\end{example}

\section{Building a Bell State}

We now assemble single- and two-qubit gates into a circuit that creates
entanglement. Prepare two qubits in $\ket{00}$, apply $H$ to qubit~$0$, then
apply $\CNOT$ with qubit~$0$ as control and qubit~$1$ as target.

\begin{theorem}[Bell-state preparation]
\label{l2:thm:bell}
With qubit~$0$ the control and qubit~$1$ the target,
\[
  \CNOT\,(H\otimes I)\,\ket{00}
    = \frac{1}{\sqrt2}\big(\ket{00}+\ket{11}\big)
    =: \ket{\Phi^{+}}.
\]
\end{theorem}

\begin{proof}
First apply $H\otimes I$. Since $H\ket{0}=\tfrac{1}{\sqrt2}(\ket{0}+\ket{1})$
and $I\ket{0}=\ket{0}$,
\[
  (H\otimes I)\ket{00}
    = \big(H\ket{0}\big)\otimes\big(I\ket{0}\big)
    = \frac{1}{\sqrt2}(\ket{0}+\ket{1})\otimes\ket{0}
    = \frac{1}{\sqrt2}\big(\ket{00}+\ket{10}\big).
\]
Now apply $\CNOT$ term by term using $\CNOT\ket{a,b}=\ket{a,a\oplus b}$. On
$\ket{00}$ the control is $0$, so the state is unchanged. On $\ket{10}$ the
control is $1$, so the target flips: $\ket{10}\mapsto\ket{11}$. Hence
\[
  \CNOT\,\frac{1}{\sqrt2}\big(\ket{00}+\ket{10}\big)
    = \frac{1}{\sqrt2}\big(\ket{00}+\ket{11}\big)=\ket{\Phi^{+}}. \qedhere
\]
\end{proof}

The state $\ket{\Phi^{+}}$ is entangled: it cannot be written as a product
$\ket{u}\otimes\ket{v}$. Indeed, a product state has amplitudes $u_{a}v_{b}$ that
factor, so the $2\times2$ amplitude matrix
$\left(\begin{smallmatrix}u_{0}v_{0}&u_{0}v_{1}\\u_{1}v_{0}&u_{1}v_{1}\end{smallmatrix}\right)$
has determinant $0$; but $\ket{\Phi^{+}}$ has amplitude matrix
$\tfrac{1}{\sqrt2}\left(\begin{smallmatrix}1&0\\0&1\end{smallmatrix}\right)$ with
determinant $\tfrac12\neq0$, a contradiction.

\begin{example}[$\CNOT$ on $\ket{+}\ket{0}$]
Since $\ket{+}=H\ket{0}$, the input $\ket{+}\ket{0}=(H\otimes I)\ket{00}$ is
exactly the intermediate state of Theorem~\ref{l2:thm:bell}. Therefore
$\CNOT\ket{+}\ket{0}=\tfrac{1}{\sqrt2}(\ket{00}+\ket{11})=\ket{\Phi^{+}}$: the
$\CNOT$ turns a product of a superposition and $\ket{0}$ into a maximally
entangled Bell pair. This is the same computation viewed with the $H$ absorbed
into the first factor.
\end{example}

\section{Universal Gate Sets and Compilation}

\begin{definition}[Universal gate set]
A finite set $\mathcal{G}$ of gates is \emph{universal} if the unitaries
generated by applying gates from $\mathcal{G}$ to arbitrary subsets of qubits are
dense in the group of all $n$-qubit unitaries (up to global phase): for every
target unitary $U$ and every $\varepsilon>0$ there is a finite circuit $V$ over
$\mathcal{G}$ with $\norm{U-V}\le\varepsilon$.
\end{definition}

Two facts underlie this definition. First, any $n$-qubit unitary can be written
\emph{exactly} as a product of single-qubit and two-qubit gates,
$U=U_{m}U_{m-1}\cdots U_{1}$, though $m$ may be exponential in $n$. Second, a
fixed \emph{finite} set already suffices to approximate every gate: the set
\[
  \mathcal{G}=\{H,\,T,\,\CNOT\}
\]
is universal. We settle for approximation because a discrete set cannot hit a
continuum of unitaries exactly, but it can come within any $\varepsilon$.

\begin{remark}[Solovay--Kitaev cost]
The Solovay--Kitaev theorem quantifies the price of this discretization: any
single-qubit gate can be approximated to error $\varepsilon$ by a circuit of
$O\!\big(\mathrm{polylog}(1/\varepsilon)\big)$ gates from a universal set, and an
$m$-gate circuit can be recompiled into the finite set using only
\[
  O\!\big(m\,\mathrm{polylog}(m/\varepsilon)\big)
\]
gates, found by an efficient classical algorithm. The overhead is thus only
polylogarithmic in the inverse precision, which is why compiling ideal circuits
to a fixed hardware gate set is practical.
\end{remark}

Running an abstract circuit on a real device requires \emph{compilation}:
(i) rewriting logical gates such as $H$, $T$, $\CNOT$ into the machine's
native gate set; and (ii) inserting \textsc{swap} gates to route qubits, because
hardware qubits are not all pairwise connected. Throughout one minimizes gate
count and depth to suppress noise and maximize fidelity.

\section{The No-Cloning Theorem}
\label{l2:sec:nocloning}

Classically, copying a bit is trivial. Quantum mechanically it is impossible to
copy an \emph{unknown} state, a fact discovered by Wootters, Zurek, and Dieks in
1982. A copier would be a gate $U$ that, given the state to copy and a fixed
blank register $\ket{s}$, outputs two copies.

\begin{theorem}[No-Cloning Theorem]
\label{l2:thm:nocloning}
There is no unitary $U$ on $\mathcal{H}\otimes\mathcal{H}$ and blank state
$\ket{s}\in\mathcal{H}$ such that
\[
  U\big(\ket{\psi}\otimes\ket{s}\big)=\ket{\psi}\otimes\ket{\psi}
  \qquad\text{for every normalized } \ket{\psi}\in\mathcal{H},
\]
unless $\dim\mathcal{H}=1$. More precisely, any such $U$ can clone two states
$\ket{x_{1}},\ket{x_{2}}$ simultaneously only if they are equal or orthogonal.
\end{theorem}

\begin{proof}
Suppose such $U$ and $\ket{s}$ exist. Apply the cloning equation to two
normalized states $\ket{x_{1}}$ and $\ket{x_{2}}$:
\[
  U\ket{x_{1}}\ket{s}=\ket{x_{1}}\ket{x_{1}},\qquad
  U\ket{x_{2}}\ket{s}=\ket{x_{2}}\ket{x_{2}}.
\]
Take the inner product of these two equations. On the left, because $U$ is
unitary it preserves inner products (Theorem~\ref{l2:thm:unitary-preserves}), and
$\braket{s|s}=1$, so
\[
  \big(\bra{x_{1}}\bra{s}\big)U^{\dagger}U\big(\ket{x_{2}}\ket{s}\big)
   = \braket{x_{1}|x_{2}}\,\braket{s|s}
   = \braket{x_{1}|x_{2}}.
\]
On the right, the inner product of $\ket{x_{1}}\ket{x_{1}}$ with
$\ket{x_{2}}\ket{x_{2}}$ factors across the tensor product,
\[
  \big(\bra{x_{1}}\bra{x_{1}}\big)\big(\ket{x_{2}}\ket{x_{2}}\big)
   = \braket{x_{1}|x_{2}}\,\braket{x_{1}|x_{2}}
   = \braket{x_{1}|x_{2}}^{2}.
\]
Equating the two sides,
\[
  \braket{x_{1}|x_{2}} = \braket{x_{1}|x_{2}}^{2}.
\]
Writing $z=\braket{x_{1}|x_{2}}$, this reads $z=z^{2}$, i.e. $z(z-1)=0$, so
$z\in\{0,1\}$. Thus $\ket{x_{1}}$ and $\ket{x_{2}}$ are either identical
($z=1$) or orthogonal ($z=0$). A universal cloner would need this for
\emph{all} pairs, but a space of dimension $\geq 2$ contains states with inner
product strictly between $0$ and $1$ (e.g. $\ket{0}$ and $\ket{+}$, with
$\braket{0|+}=\tfrac{1}{\sqrt2}$). This contradiction proves no universal cloner
exists; only mutually orthogonal states can be cloned.
\end{proof}

The essence of the proof is that cloning
$\ket{\psi}\mapsto\ket{\psi}\otimes\ket{\psi}$ is a nonlinear map of $\ket{\psi}$
(it is quadratic in the amplitudes), whereas every quantum gate is linear;
linearity forces the inner-product constraint $z=z^{2}$, which no continuum of
states can satisfy.

\begin{remark}[Consequences]\leavevmode
\begin{itemize}[nosep]
  \item \emph{No broadcasting.} Since an unknown state cannot even be duplicated,
        it cannot be broadcast to several parties.
  \item \emph{$\CNOT$ does not clone superpositions.} For
        $\ket{x}=a\ket{0}+b\ket{1}$,
        \[
          \CNOT\ket{x}\ket{0}=a\ket{00}+b\ket{11}\neq\ket{x}\ket{x}
          =\ket{x}\otimes(a\ket{0}+b\ket{1}),
        \]
        because the right side contains cross terms $ab\,\ket{01},\,ab\,\ket{10}$
        absent on the left. $\CNOT$ copies only computational-basis
        (classical) information, consistent with Theorem~\ref{l2:thm:nocloning}.
  \item \emph{Measurement destroys unknown states.} A projective measurement
        sends $\ket{\psi}\mapsto P_{m}\ket{\psi}/\sqrt{p_{m}}$ with
        $p_{m}=\bra{\psi}P_{m}\ket{\psi}$; the outcome is random and the
        pre-measurement superposition is irreversibly lost.
  \item \emph{BB84 security.} Because an eavesdropper can neither copy the
        transmitted qubits nor measure them without disturbing non-orthogonal
        states, interception in the BB84 protocol necessarily introduces
        detectable errors, giving quantum key distribution its security.
\end{itemize}
\end{remark}

\begin{theorem}[No-Deleting Theorem]
\label{l2:thm:nodeleting}
There is no unitary $U$ that deletes an unknown state against a copy, i.e. no
$U$ and fixed target $\ket{\Sigma}$ with
\[
  U\big(\ket{\psi}\otimes\ket{\psi}\big)=\ket{\psi}\otimes\ket{\Sigma}
  \qquad\text{for all normalized } \ket{\psi},
\]
on a Hilbert space of dimension $\geq 2$.
\end{theorem}

\begin{proof}
Assume such $U$ exists and apply it to two states $\ket{x_{1}},\ket{x_{2}}$:
\[
  U\ket{x_{1}}\ket{x_{1}}=\ket{x_{1}}\ket{\Sigma},\qquad
  U\ket{x_{2}}\ket{x_{2}}=\ket{x_{2}}\ket{\Sigma}.
\]
Taking inner products and using unitarity on the left and factoring on the right,
\[
  \braket{x_{1}|x_{2}}^{2}=\braket{x_{1}|x_{2}}\,\braket{\Sigma|\Sigma}
                          =\braket{x_{1}|x_{2}},
\]
since $\braket{\Sigma|\Sigma}=1$. Again $z=\braket{x_{1}|x_{2}}$ satisfies
$z^{2}=z$, so $z\in\{0,1\}$. As in Theorem~\ref{l2:thm:nocloning}, this fails for
non-orthogonal states, so no universal deleting operation exists.
\end{proof}

\section{Problems}

\begin{exercise}\label{l2:ex:hth}
Verify by matrix multiplication that $H$ is unitary and self-inverse
($H^{\dagger}=H$ and $H^{2}=I$). Then, using Lemma~\ref{l2:lem:conj}, compute
$HYH$.
\end{exercise}

\begin{exercise}\label{l2:ex:st}
Show that $S=T^{2}$ and $Z=S^{2}$, and conclude $Z=T^{4}$. What single-qubit
rotation angle about the $z$-axis does $T$ implement?
\end{exercise}

\begin{exercise}\label{l2:ex:cnot-conj}
Prove that $(I\otimes H)\,\CNOT\,(I\otimes H)$ acts as a controlled-$Z$ gate,
i.e. it applies a phase $-1$ exactly to $\ket{11}$ and fixes the other basis
states. (Here $H$ acts on the target qubit and the control is qubit~$0$.)
\end{exercise}

\begin{exercise}\label{l2:ex:reverse}
A single-qubit circuit applies, in order, $T$ then $H$ then $T$. Write its
overall unitary $U$ as a matrix product in the correct order and evaluate it.
\end{exercise}

\begin{exercise}\label{l2:ex:cnot-plus}
Compute $\CNOT\ket{+}\ket{+}$ and decide whether the result is entangled.
\end{exercise}

\begin{exercise}\label{l2:ex:noclone-orth}
Show explicitly that the orthogonal states $\ket{0},\ket{1}$ \emph{can} be cloned
by a single gate: exhibit a unitary $U$ and blank $\ket{s}$ with
$U\ket{0}\ket{s}=\ket{0}\ket{0}$ and $U\ket{1}\ket{s}=\ket{1}\ket{1}$. Why does
this not contradict Theorem~\ref{l2:thm:nocloning}?
\end{exercise}

\begin{exercise}\label{l2:ex:swap}
Show that three alternating $\CNOT$s implement \textsc{swap}. Concretely, with
$C_{01}$ the $\CNOT$ with control $0$/target $1$ and $C_{10}$ the reverse, verify
$C_{01}C_{10}C_{01}\ket{a,b}=\ket{b,a}$ for all $a,b\in\{0,1\}$.
\end{exercise}

\subsection*{Solutions}

\begin{solution}
\textbf{(Exercise~\ref{l2:ex:hth}.)} $H$ has real symmetric entries, so
$H^{\dagger}=H^{\top}=H$. Then
$H^{2}=\tfrac12\left(\begin{smallmatrix}1&1\\1&-1\end{smallmatrix}\right)
\left(\begin{smallmatrix}1&1\\1&-1\end{smallmatrix}\right)
=\tfrac12\left(\begin{smallmatrix}2&0\\0&2\end{smallmatrix}\right)=I$, which shows
$H^{\dagger}H=H^{2}=I$, so $H$ is unitary and self-inverse. For $HYH$, use
$Y=iXZ$ (Proposition~\ref{l2:prop:yxz}) and $H^{2}=I$:
\[
  HYH=iH X Z H = i(HXH)(HZH)=i\,Z\,X
     = i\begin{pmatrix}1&0\\0&-1\end{pmatrix}
       \begin{pmatrix}0&1\\1&0\end{pmatrix}
     = i\begin{pmatrix}0&1\\-1&0\end{pmatrix}
     = \begin{pmatrix}0&i\\-i&0\end{pmatrix}=-Y.
\]
So $HYH=-Y$.
\end{solution}

\begin{solution}
\textbf{(Exercise~\ref{l2:ex:st}.)} All are diagonal. $T^{2}
=\operatorname{diag}(1,e^{i\pi/4})^{2}=\operatorname{diag}(1,e^{i\pi/2})
=\operatorname{diag}(1,i)=S$. Next $S^{2}=\operatorname{diag}(1,i)^{2}
=\operatorname{diag}(1,-1)=Z$. Combining, $Z=S^{2}=(T^{2})^{2}=T^{4}$. Since
$T=\operatorname{diag}(1,e^{i\pi/4})$ multiplies the $\ket{1}$ amplitude by
$e^{i\pi/4}$, it implements a $z$-axis rotation by angle $\pi/4$ in the relative
phase (a $\pi/8$ rotation of the Bloch vector about the $z$-axis, whence the
common name ``$\pi/8$ gate'').
\end{solution}

\begin{solution}
\textbf{(Exercise~\ref{l2:ex:cnot-conj}.)} Let $V=(I\otimes H)\,\CNOT\,(I\otimes H)$.
Act on a basis state $\ket{a}\ket{b}$. On the target, the first $H$ gives
$\ket{a}\,H\ket{b}$; $\CNOT$ maps this to $\ket{a}\,X^{a}H\ket{b}$; the final $H$
gives $\ket{a}\,H X^{a} H\ket{b}$. By Lemma~\ref{l2:lem:conj}, $HX H=Z$ and for
$a=0$ we get $H^{2}=I$. Thus $HX^{a}H=Z^{a}$, so
$V\ket{a}\ket{b}=\ket{a}\,Z^{a}\ket{b}$. This applies $Z$ to the target exactly
when $a=1$; and $Z\ket{1}=-\ket{1}$, so the only state acquiring the phase $-1$
is $\ket{11}$, with $\ket{00},\ket{01},\ket{10}$ fixed. Hence $V$ is the
controlled-$Z$ gate, which is symmetric in the two qubits.
\end{solution}

\begin{solution}
\textbf{(Exercise~\ref{l2:ex:reverse}.)} The circuit applies $T$, then $H$, then
$T$; by the reverse-order rule the unitary is $U=T\,H\,T$. Compute
$HT=\tfrac{1}{\sqrt2}\left(\begin{smallmatrix}1&1\\1&-1\end{smallmatrix}\right)
\left(\begin{smallmatrix}1&0\\0&e^{i\pi/4}\end{smallmatrix}\right)
=\tfrac{1}{\sqrt2}\left(\begin{smallmatrix}1&e^{i\pi/4}\\1&-e^{i\pi/4}\end{smallmatrix}\right)$,
then
\[
  U=T(HT)=\begin{pmatrix}1&0\\0&e^{i\pi/4}\end{pmatrix}
     \frac{1}{\sqrt2}\begin{pmatrix}1&e^{i\pi/4}\\1&-e^{i\pi/4}\end{pmatrix}
   =\frac{1}{\sqrt2}\begin{pmatrix}1&e^{i\pi/4}\\ e^{i\pi/4}&-e^{i\pi/2}\end{pmatrix}
   =\frac{1}{\sqrt2}\begin{pmatrix}1&e^{i\pi/4}\\ e^{i\pi/4}&-i\end{pmatrix}.
\]
One verifies $U^{\dagger}U=I$, as it must for a product of gates.
\end{solution}

\begin{solution}
\textbf{(Exercise~\ref{l2:ex:cnot-plus}.)} Write
$\ket{+}\ket{+}=\tfrac12(\ket{00}+\ket{01}+\ket{10}+\ket{11})$. Apply
$\CNOT$ (control $0$) term by term: $\ket{00}\to\ket{00}$,
$\ket{01}\to\ket{01}$, $\ket{10}\to\ket{11}$, $\ket{11}\to\ket{10}$. The sum is
invariant, so
$\CNOT\ket{+}\ket{+}=\tfrac12(\ket{00}+\ket{01}+\ket{11}+\ket{10})=\ket{+}\ket{+}$.
The result is a product state, hence \emph{not} entangled: $\ket{+}\ket{+}$ is a
$+1$ eigenstate of this $\CNOT$.
\end{solution}

\begin{solution}
\textbf{(Exercise~\ref{l2:ex:noclone-orth}.)} Take $\ket{s}=\ket{0}$ and let
$U=\CNOT$ with control $0$, target $1$. Then $\CNOT\ket{0}\ket{0}=\ket{0}\ket{0}$
and $\CNOT\ket{1}\ket{0}=\ket{1}\ket{1}$, so both orthogonal basis states are
copied. This does not contradict Theorem~\ref{l2:thm:nocloning}: the theorem only
forbids a \emph{single} gate that copies \emph{every} state. This $\CNOT$
copies $\ket{0},\ket{1}$ but, as shown in the No-Cloning remark, fails to copy a
superposition $a\ket{0}+b\ket{1}$, instead producing the entangled
$a\ket{00}+b\ket{11}$. Cloning a fixed orthonormal set is exactly the ``classical
copy'' the theorem permits.
\end{solution}

\begin{solution}
\textbf{(Exercise~\ref{l2:ex:swap}.)} Track the pair $(a,b)$ under the three gates,
recalling $C_{01}:(a,b)\mapsto(a,a\oplus b)$ and $C_{10}:(a,b)\mapsto
(a\oplus b,b)$. Starting from $(a,b)$:
\[
  (a,b)\ \xrightarrow{C_{01}}\ (a,\,a\oplus b)
       \ \xrightarrow{C_{10}}\ (a\oplus(a\oplus b),\,a\oplus b)=(b,\,a\oplus b)
       \ \xrightarrow{C_{01}}\ (b,\,b\oplus(a\oplus b))=(b,a).
\]
Since $a\oplus a=0$ and $b\oplus b=0$, the net effect is $(a,b)\mapsto(b,a)$ for
all $a,b$, which is exactly \textsc{swap}. As each $\CNOT$ is a permutation of
basis states, so is the composite, confirming
$\mathrm{SWAP}\ket{a,b}=\ket{b,a}$.
\end{solution}

\chapter{The Deutsch--Jozsa and Bernstein--Vazirani Algorithms}

The algorithms in this chapter are historically the first to demonstrate, in a
clean and \emph{deterministic} way, that a quantum computer can outperform any
classical computer. The tasks they solve are artificial and of little practical
value, yet the tools they introduce---the oracle model, phase kickback, and
Hadamard interference---form the conceptual backbone of nearly every later
quantum algorithm, from Simon's problem to Shor's factoring algorithm. Master
these three ideas and the rest of the course rests on solid ground.

Throughout, $f$ denotes an unknown Boolean function that we may access only
through an oracle. Our measure of cost is the \emph{number of oracle queries},
not the number of elementary gates. We will see that a task requiring
exponentially many classical queries can be solved with a single quantum query.

\section{The Query and Oracle Model}

We are given a Boolean function $f:\{0,1\}^n\to\{0,1\}$ whose internal structure
is hidden from us. We do not have the code of $f$; we may only feed inputs and
receive outputs through a black box. To use $f$ inside a quantum circuit we need
its action to be reversible, since every quantum gate is unitary. The standard
device for this is the following.

\begin{definition}[Oracle / query model]
\label{l3:def:oracle}
Let $f:\{0,1\}^n\to\{0,1\}$. The \emph{oracle} (or \emph{quantum oracle}) for $f$
is the operator $\Uf$ acting on $n+1$ qubits---an $n$-qubit \emph{input register}
and a single-qubit \emph{target register}---defined on the computational basis by
\[
  \Uf\ket{x}\ket{y}=\ket{x}\ket{y\oplus f(x)},
  \qquad x\in\{0,1\}^n,\ y\in\{0,1\},
\]
where $\oplus$ denotes addition modulo $2$. The \emph{query complexity} of a task
is the minimum number of applications of $\Uf$ (over all algorithms that succeed)
needed to solve it.
\end{definition}

The exclusive-or in the target register---rather than simply overwriting $y$ with
$f(x)$---is precisely what makes the oracle reversible.

\begin{lemma}[The oracle is unitary]
\label{l3:lem:unitary}
For every $f:\{0,1\}^n\to\{0,1\}$, the operator $\Uf$ of Definition~\ref{l3:def:oracle}
is unitary and is its own inverse: $\Uf^{\dagger}\Uf=\Uf^2=\id$.
\end{lemma}

\begin{proof}
The set $\{\ket{x}\ket{y}\}$ over $x\in\{0,1\}^n$, $y\in\{0,1\}$ is an orthonormal
basis of the $2^{n+1}$-dimensional state space. The map
$(x,y)\mapsto(x,\,y\oplus f(x))$ is a bijection of the index set to itself, because
it is inverted by itself: applying it twice gives
$(x,\,y\oplus f(x)\oplus f(x))=(x,y)$ since $f(x)\oplus f(x)=0$. Hence $\Uf$ merely
permutes an orthonormal basis, so it is unitary, and $\Uf^2=\id$. Equivalently, in
matrix form $\braket{x'y'|\Uf^{\dagger}\Uf|xy}=\delta_{xx'}\delta_{yy'}$, i.e.
$\Uf^{\dagger}\Uf=\id$.
\end{proof}

\paragraph{Classical versus quantum access.}
A classical algorithm can only submit one input $x$ at a time and read back the
single bit $f(x)$; setting the target to $y=0$ gives $\Uf\ket{x}\ket{0}=\ket{x}\ket{f(x)}$.
To learn a \emph{global} property of $f$ this way may force us to examine almost
all $2^n$ inputs. A quantum algorithm can instead prepare a uniform superposition
of all inputs and apply $\Uf$ once, acting on every input in parallel:
\[
  \Uf\Big(\tfrac{1}{\sqrt{N}}\sum_{x\in\{0,1\}^n}\ket{x}\Big)\ket{0}
  =\tfrac{1}{\sqrt{N}}\sum_{x}\ket{x}\ket{f(x)},
  \qquad N=2^n.
\]
Parallelism alone is not enough---measuring this state returns one random pair
$(x,f(x))$, no better than a classical query. The extra ingredient that converts
parallelism into advantage is \emph{interference}, engineered through phase
kickback and the Hadamard transform.

\section{Phase Kickback}

The oracle writes $f(x)$ into the target register as a bit. We would prefer to
have $f(x)$ appear as a \emph{phase} on the input register, since phases are what
interfere. The following observation---perhaps the single most important trick in
the chapter---accomplishes exactly this by choosing the target register cleverly.

\begin{lemma}[Phase kickback]
\label{l3:lem:kickback}
Let $\ket{-}=\tfrac{1}{\sqrt2}(\ket{0}-\ket{1})$. Then for every $x\in\{0,1\}^n$,
\[
  \Uf\ket{x}\ket{-}=(-1)^{f(x)}\ket{x}\ket{-}.
\]
In particular the target qubit is left in $\ket{-}$ and may be discarded, so on
the input register the oracle acts as the effective diagonal unitary
$\tilde U_f\ket{x}=(-1)^{f(x)}\ket{x}$.
\end{lemma}

\begin{proof}
Expand the target register and apply Definition~\ref{l3:def:oracle} to each branch:
\[
  \Uf\ket{x}\ket{-}
  =\tfrac{1}{\sqrt2}\big(\Uf\ket{x}\ket{0}-\Uf\ket{x}\ket{1}\big)
  =\tfrac{1}{\sqrt2}\big(\ket{x}\ket{f(x)}-\ket{x}\ket{1\oplus f(x)}\big).
\]
Consider the two cases for the bit $f(x)$. If $f(x)=0$ the target factor is
$\tfrac{1}{\sqrt2}(\ket{0}-\ket{1})=\ket{-}=(+1)\ket{-}$. If $f(x)=1$ it is
$\tfrac{1}{\sqrt2}(\ket{1}-\ket{0})=-\ket{-}=(-1)\ket{-}$. Both cases are captured
by the single formula $(-1)^{f(x)}\ket{-}$, whence
$\Uf\ket{x}\ket{-}=(-1)^{f(x)}\ket{x}\ket{-}$. Since the target factor is unchanged
and unentangled from the input register, it factors out and can be dropped.
\end{proof}

\begin{remark}
The name \emph{kickback} reflects that the eigenvalue $(-1)^{f(x)}$ of the target
qubit's controlled action is ``kicked back'' onto the control (input) register.
The target qubit acts purely as a passive signal state. This is the same mechanism
that, with a target eigenstate of a general unitary, powers quantum phase
estimation later in the course.
\end{remark}

Applying $\Uf$ to a uniform superposition with the target in $\ket{-}$ therefore
produces, on the input register,
\[
  \tfrac{1}{\sqrt{N}}\sum_{x}\ket{x}\;\xmapsto{\ \Uf\ }\;
  \tfrac{1}{\sqrt{N}}\sum_{x}(-1)^{f(x)}\ket{x}.
\]
Every value of $f$ is now recorded in a relative sign. It remains to read out
these signs; the tool for that is the Hadamard transform.

\section{The Hadamard Transform}

We will use repeatedly the action of $H^{\otimes n}$ on a basis state. Recall for a
single qubit $H\ket{a}=\tfrac{1}{\sqrt2}\sum_{b\in\{0,1\}}(-1)^{ab}\ket{b}$.

\begin{lemma}[Hadamard identity]
\label{l3:lem:had}
For $x\in\{0,1\}^n$, writing $x\cdot z=\bigoplus_{i=1}^{n}x_i z_i$ for the bitwise
inner product modulo $2$,
\[
  H^{\otimes n}\ket{x}=\frac{1}{\sqrt{2^{\,n}}}\sum_{z\in\{0,1\}^n}(-1)^{x\cdot z}\ket{z}.
\]
\end{lemma}

\begin{proof}
Take the tensor product of the single-qubit rule over the $n$ qubits:
\[
  H^{\otimes n}\ket{x}
  =\bigotimes_{i=1}^{n}\Big(\tfrac{1}{\sqrt2}\sum_{z_i\in\{0,1\}}(-1)^{x_i z_i}\ket{z_i}\Big)
  =\frac{1}{\sqrt{2^{\,n}}}\sum_{z\in\{0,1\}^n}(-1)^{\sum_i x_i z_i}\ket{z}.
\]
Because the phase is $\pm1$, the exponent $\sum_i x_i z_i$ may be reduced modulo
$2$, giving $(-1)^{x\cdot z}$.
\end{proof}

We also record the orthogonality relation that makes the interference exact:
\[
  \frac{1}{2^{\,n}}\sum_{x\in\{0,1\}^n}(-1)^{x\cdot w}
  =\delta_{w,0^n},
  \qquad w\in\{0,1\}^n,
\]
which holds because the sum factors as $\prod_{i=1}^{n}\tfrac12\sum_{x_i}(-1)^{x_i w_i}$,
and each factor equals $1$ if $w_i=0$ and $0$ if $w_i=1$.

\section{The Deutsch Algorithm}

Consider the simplest nontrivial case, $n=1$, so $f:\{0,1\}\to\{0,1\}$. There are
four such functions. We wish to decide a global property of $f$ using as few
queries as possible.

\begin{definition}[Constant and balanced functions]
\label{l3:def:constbal}
A function $f:\{0,1\}^n\to\{0,1\}$ is \emph{constant} if $f(x)$ takes the same value
on all $2^n$ inputs. It is \emph{balanced} if $f(x)=0$ on exactly half of the inputs
and $f(x)=1$ on the other half. For $n=1$, $f$ is constant iff $f(0)=f(1)$ and
balanced iff $f(0)\neq f(1)$.
\end{definition}

Classically, deciding whether $f$ is constant or balanced for $n=1$ requires
\emph{two} queries: knowing $f(0)$ alone tells us nothing about $f(1)$. Deutsch's
algorithm decides it with a single query.

\begin{theorem}[Correctness of the Deutsch algorithm]
\label{l3:thm:deutsch}
Prepare the two-qubit state $\ket{0}\ket{1}$, apply $H$ to each qubit, apply $\Uf$,
apply $H$ to the first qubit, and measure the first qubit. The outcome is $0$ if
$f$ is constant and $1$ if $f$ is balanced. Exactly one query to $\Uf$ is used.
\end{theorem}

\begin{proof}
After the first Hadamards,
$H\ket{0}\otimes H\ket{1}=\ket{+}\ket{-}=\tfrac{1}{\sqrt2}(\ket{0}+\ket{1})\ket{-}$.
Applying $\Uf$ and using phase kickback (Lemma~\ref{l3:lem:kickback}) on each branch,
\[
  \Uf\,\tfrac{1}{\sqrt2}\big(\ket{0}+\ket{1}\big)\ket{-}
  =\tfrac{1}{\sqrt2}\big((-1)^{f(0)}\ket{0}+(-1)^{f(1)}\ket{1}\big)\ket{-}.
\]
Now apply $H$ to the first qubit, using $H\ket{0}=\tfrac{1}{\sqrt2}(\ket0+\ket1)$
and $H\ket{1}=\tfrac{1}{\sqrt2}(\ket0-\ket1)$:
\[
  \tfrac{1}{2}\Big[\big((-1)^{f(0)}+(-1)^{f(1)}\big)\ket{0}
  +\big((-1)^{f(0)}-(-1)^{f(1)}\big)\ket{1}\Big]\ket{-}.
\]
If $f$ is constant, $f(0)=f(1)$, the $\ket{1}$ amplitude vanishes and the
$\ket{0}$ amplitude is $\pm1$; the first qubit is $\pm\ket{0}$ and measurement
yields $0$ with certainty. If $f$ is balanced, $f(0)\neq f(1)$, the $\ket{0}$
amplitude vanishes and the first qubit is $\pm\ket{1}$; measurement yields $1$
with certainty. The oracle is called once.
\end{proof}

Notice that the answer is read from the \emph{first} qubit, even though classically
one would expect the ``information'' to live in the target qubit that received
$f(x)$. This inversion is the hallmark of interference.

\begin{example}[Deutsch on $f(a)=a\oplus 1$]
\label{l3:ex:deutsch}
Take $f(0)=1$, $f(1)=0$, which is balanced. After the query the first register is
$\tfrac{1}{\sqrt2}\big((-1)^{1}\ket0+(-1)^{0}\ket1\big)=\tfrac{1}{\sqrt2}(-\ket0+\ket1)=-\ket{-}$.
Applying $H$ gives $H(-\ket{-})=-\ket{1}$, so the measurement returns $1$,
correctly reporting ``balanced.'' By contrast, for the constant function
$f\equiv 1$ the first register after the query is
$\tfrac{1}{\sqrt2}(-\ket0-\ket1)=-\ket{+}$, and $H(-\ket{+})=-\ket0$ returns $0$,
reporting ``constant.''
\end{example}

\section{The Deutsch--Jozsa Algorithm}

We now generalize to $n$ qubits under the following \emph{promise}: $f$ is either
constant or balanced (Definition~\ref{l3:def:constbal}), and we are guaranteed that
one of these holds. The task is to decide which. Classically, in the worst case
one must query $2^{n-1}+1$ inputs: after seeing $2^{n-1}$ identical outputs the
function could still be either constant or balanced, and one more query is needed.
Deutsch--Jozsa decides it with a single quantum query.

The circuit is the unifying pattern $H^{\otimes n}$--$\Uf$--$H^{\otimes n}$ on the
$n$-qubit input register, with the target qubit held in $\ket{-}$.

\begin{theorem}[Correctness of Deutsch--Jozsa]
\label{l3:thm:dj}
Prepare $\ket{0}^{\otimes n}$ on the input register and $\ket{-}$ on the target.
Apply $H^{\otimes n}$, then $\Uf$, then $H^{\otimes n}$, and measure the input
register. If $f$ is constant the outcome is $0^n$ with probability $1$; if $f$ is
balanced the outcome $0^n$ occurs with probability $0$. Thus measuring all zeros
certifies ``constant,'' and any nonzero outcome certifies ``balanced.'' One query
to $\Uf$ is used.
\end{theorem}

\begin{proof}
By Lemma~\ref{l3:lem:had}, $H^{\otimes n}\ket{0}^{\otimes n}
=\tfrac{1}{\sqrt N}\sum_{x}\ket{x}$ with $N=2^n$. Applying $\Uf$ with the target
in $\ket{-}$ triggers phase kickback (Lemma~\ref{l3:lem:kickback}) on every branch,
giving (dropping the unchanged target factor)
\[
  \tfrac{1}{\sqrt N}\sum_{x}(-1)^{f(x)}\ket{x}.
\]
Apply $H^{\otimes n}$ once more and use Lemma~\ref{l3:lem:had} again:
\[
  H^{\otimes n}\tfrac{1}{\sqrt N}\sum_{x}(-1)^{f(x)}\ket{x}
  =\frac{1}{N}\sum_{z}\Big(\sum_{x}(-1)^{f(x)}(-1)^{x\cdot z}\Big)\ket{z}.
\]
The amplitude of the outcome $z=0^n$ is obtained by setting $z=0^n$, so $x\cdot z=0$:
\[
  \alpha_{0^n}=\frac{1}{N}\sum_{x\in\{0,1\}^n}(-1)^{f(x)}.
\]
If $f$ is constant, every term is the same sign and $\alpha_{0^n}=\pm1$, so the
outcome is $0^n$ with probability $\abs{\alpha_{0^n}}^2=1$. If $f$ is balanced,
exactly $N/2$ terms equal $+1$ and $N/2$ equal $-1$, so they cancel and
$\alpha_{0^n}=0$: the outcome $0^n$ never occurs. In either case a single query
suffices, and the two cases are distinguished with certainty.
\end{proof}

\section{The Bernstein--Vazirani Algorithm}

The Deutsch--Jozsa algorithm returns a single bit of information (constant vs.\
balanced). The Bernstein--Vazirani algorithm uses the very same circuit to extract
an entire hidden $n$-bit string in one query, provided $f$ has a special linear
structure.

\begin{definition}[Bernstein--Vazirani function]
\label{l3:def:bv}
Fix a hidden string $s\in\{0,1\}^n$. The associated \emph{linear} Boolean function is
\[
  f_s(x)=s\cdot x=\bigoplus_{i=1}^{n}s_i x_i \pmod 2,
\]
the parity of the number of positions at which $s$ and $x$ are both $1$. The
Bernstein--Vazirani problem is: given oracle access to $\Uf$ for such an $f_s$,
determine $s$.
\end{definition}

Classically, one recovers $s$ bit by bit. Querying $x=e_i$ (the string with a
single $1$ in position $i$) returns $f_s(e_i)=s_i$; this determines one bit per
query, so $n$ queries are needed and sufficient. The quantum algorithm needs one.

\begin{theorem}[Correctness of Bernstein--Vazirani]
\label{l3:thm:bv}
For $f=f_s$ of Definition~\ref{l3:def:bv}, running the circuit of
Theorem~\ref{l3:thm:dj}---$H^{\otimes n}$, then $\Uf$, then $H^{\otimes n}$, with the
target in $\ket{-}$---leaves the input register in the state $\ket{s}$. Measuring
recovers $s$ with certainty using a single query.
\end{theorem}

\begin{proof}
As before, after $H^{\otimes n}$ and the single query with phase kickback the input
register is
\[
  \tfrac{1}{\sqrt N}\sum_{x}(-1)^{f_s(x)}\ket{x}
  =\tfrac{1}{\sqrt N}\sum_{x}(-1)^{s\cdot x}\ket{x}.
\]
Applying $H^{\otimes n}$ and using Lemma~\ref{l3:lem:had},
\[
  H^{\otimes n}\tfrac{1}{\sqrt N}\sum_{x}(-1)^{s\cdot x}\ket{x}
  =\frac{1}{N}\sum_{z}\Big(\sum_{x}(-1)^{s\cdot x}(-1)^{x\cdot z}\Big)\ket{z}
  =\frac{1}{N}\sum_{z}\Big(\sum_{x}(-1)^{x\cdot(s\oplus z)}\Big)\ket{z},
\]
where $s\cdot x\oplus x\cdot z=x\cdot(s\oplus z)$ because the inner product is
linear in its second argument modulo $2$. By the orthogonality relation of
Section~3, $\tfrac{1}{N}\sum_{x}(-1)^{x\cdot(s\oplus z)}=\delta_{s\oplus z,\,0^n}
=\delta_{z,s}$. Hence only the term $z=s$ survives and the register is exactly
$\ket{s}$. Measuring in the computational basis returns $s$ with probability $1$.
\end{proof}

Equivalently, $H^{\otimes n}\ket{s}=\tfrac{1}{\sqrt N}\sum_x(-1)^{s\cdot x}\ket{x}$
by Lemma~\ref{l3:lem:had}, so the post-query state is $H^{\otimes n}\ket{s}$; a second
$H^{\otimes n}$ inverts it (since $H^{\otimes n}$ is its own inverse) and returns
$\ket{s}$.

\begin{example}[A small Bernstein--Vazirani instance]
\label{l3:ex:bv}
Let $n=3$ and $s=101$, so $f_s(x)=x_1\oplus x_3$. After $H^{\otimes 3}$ and one
query, the input register is $\tfrac{1}{\sqrt8}\sum_{x}(-1)^{x_1\oplus x_3}\ket{x}$.
Applying $H^{\otimes 3}$ collapses this to $\ket{101}$, and the measurement reads
off $s=101$ in one shot. A classical solver would instead query
$x=100$, $010$, $001$ to obtain $s_1=1$, $s_2=0$, $s_3=1$ respectively---three
queries for the same answer.
\end{example}

\section{Discussion and Significance}

\begin{remark}[Deterministic advantage and its limits]
\label{l3:rem:significance}
Several points deserve emphasis.
\begin{itemize}
\item \emph{Deterministic vs.\ bounded-error.} All three algorithms succeed with
probability exactly $1$; the advantage is not a matter of a lucky measurement. For
Deutsch--Jozsa the separation is $1$ quantum query versus $2^{n-1}+1$ classical
queries in the worst case (deterministically). If one allows a small classical
error probability, a handful of random queries distinguishes constant from
balanced with high confidence, so the \emph{bounded-error} classical query
complexity of Deutsch--Jozsa is only $O(1)$; the exponential gap there is specific
to the deterministic setting. For Bernstein--Vazirani the separation is $1$ versus
$n$ queries and survives even in the bounded-error model.
\item \emph{The unifying pattern.} Every algorithm here is the circuit
$H^{\otimes n}\,\Uf\,H^{\otimes n}$ applied to $\ket{0}^{\otimes n}$ with the target
in $\ket{-}$. Phase kickback writes $f$ into signs; the two Hadamard layers read
those signs out by interference. This template recurs throughout quantum
algorithms.
\item \emph{Query advantage $\neq$ gate-complexity advantage.} We count calls to
$\Uf$, treating the oracle as a unit-cost black box. The number of elementary
gates needed to \emph{build} $\Uf$ for a given $f$ must be analyzed separately, and
for a generic $f$ it can be large. The dramatic query separations above do not by
themselves imply a separation in total running time.
\item \emph{Relation to Simon and Shor.} Bernstein--Vazirani generalizes to
Simon's problem, where $f$ obeys the hidden-structure promise
$f(x)=f(x\oplus s)$ (equivalently $f(x)=f(y)$ iff $x\oplus y\in\{0^n,s\}$); the
same $H^{\otimes n}\,\Uf\,H^{\otimes n}$ pattern yields linear equations
$z\cdot s=0$ constraining $s$, and Simon's exponential separation holds even with
bounded error. Simon's algorithm in turn inspired Shor's factoring algorithm,
where the final Hadamard transform is replaced by the quantum Fourier transform of
a later chapter.
\end{itemize}
\end{remark}

\section{Problems}

\begin{exercise}
\label{l3:ex:p1}
List all four functions $f:\{0,1\}\to\{0,1\}$ and classify each as constant or
balanced. For each, compute by hand the first-qubit measurement outcome produced
by the Deutsch algorithm of Theorem~\ref{l3:thm:deutsch}.
\end{exercise}

\begin{exercise}
\label{l3:ex:p2}
Prove that $\Uf$ of Definition~\ref{l3:def:oracle} is Hermitian, and use this together
with $\Uf^2=\id$ to conclude that its eigenvalues are $\pm1$.
\end{exercise}

\begin{exercise}
\label{l3:ex:p3}
Show directly, without invoking Lemma~\ref{l3:lem:kickback}, that if the target qubit
is initialized in $\ket{+}=\tfrac{1}{\sqrt2}(\ket0+\ket1)$ instead of $\ket{-}$,
then $\Uf\ket{x}\ket{+}=\ket{x}\ket{+}$ for all $x$, so no phase is kicked back.
Explain why $\ket{-}$ is the useful choice.
\end{exercise}

\begin{exercise}
\label{l3:ex:p4}
For $n=2$, let $f(x_1x_2)=x_1\oplus x_2$. Verify that $f$ is balanced and trace the
Deutsch--Jozsa state through $H^{\otimes 2}$, $\Uf$, $H^{\otimes 2}$, confirming
that the amplitude of $\ket{00}$ is $0$.
\end{exercise}

\begin{exercise}
\label{l3:ex:p5}
In the Bernstein--Vazirani setting with hidden string $s$, suppose the oracle is
for $g(x)=s\cdot x\oplus c$ with an unknown constant bit $c$ (an ``affine''
function). Show that the algorithm of Theorem~\ref{l3:thm:bv} still outputs $s$ with
certainty, and identify where the constant $c$ goes.
\end{exercise}

\begin{exercise}
\label{l3:ex:p6}
Prove the orthogonality relation $\tfrac{1}{2^n}\sum_{x\in\{0,1\}^n}(-1)^{x\cdot w}
=\delta_{w,0^n}$ used in the proof of Theorem~\ref{l3:thm:bv}, by factoring the sum
over the $n$ bits.
\end{exercise}

\begin{exercise}
\label{l3:ex:p7}
Give the exact worst-case classical (deterministic) query complexity of the
Deutsch--Jozsa problem on $n$ bits and justify it. Contrast this with the
observation that a randomized classical algorithm distinguishes constant from
balanced with error at most $1/3$ using only a constant number of queries.
\end{exercise}

\section{Solutions}

\begin{solution}[Exercise~\ref{l3:ex:p1}]
The four functions are: $f\equiv0$ (constant), $f\equiv1$ (constant),
$f(a)=a$ (balanced, $f(0)=0,f(1)=1$), and $f(a)=a\oplus1$ (balanced,
$f(0)=1,f(1)=0$). By Theorem~\ref{l3:thm:deutsch} the outcome is $0$ for the two
constant functions and $1$ for the two balanced ones. Concretely, for $f\equiv0$
the post-query first register is $\ket{+}\to H\ket{+}=\ket0$; for $f\equiv1$ it is
$-\ket{+}\to-\ket0$ (outcome $0$); for $f(a)=a$ it is
$\tfrac{1}{\sqrt2}(\ket0-\ket1)=\ket{-}\to\ket1$; for $f(a)=a\oplus1$ it is
$-\ket{-}\to-\ket1$ (outcome $1$).
\end{solution}

\begin{solution}[Exercise~\ref{l3:ex:p2}]
In the computational basis $\Uf$ is a real permutation matrix
($\ket{x}\ket{y}\mapsto\ket{x}\ket{y\oplus f(x)}$). Because the permutation is an
involution (Lemma~\ref{l3:lem:unitary}), the matrix equals its own transpose, hence is
symmetric; being real and symmetric it is Hermitian, $\Uf^{\dagger}=\Uf$.
Alternatively, from $\Uf^{\dagger}\Uf=\id$ (unitary) and $\Uf^2=\id$ (involution) we
get $\Uf^{\dagger}=\Uf^{-1}=\Uf$. If $\Uf\ket{\psi}=\lambda\ket{\psi}$ then
$\ket{\psi}=\Uf^2\ket{\psi}=\lambda^2\ket{\psi}$, so $\lambda^2=1$ and
$\lambda\in\{+1,-1\}$.
\end{solution}

\begin{solution}[Exercise~\ref{l3:ex:p3}]
Expanding as in Lemma~\ref{l3:lem:kickback},
\[
  \Uf\ket{x}\ket{+}
  =\tfrac{1}{\sqrt2}\big(\ket{x}\ket{f(x)}+\ket{x}\ket{1\oplus f(x)}\big)
  =\tfrac{1}{\sqrt2}\ket{x}\big(\ket{f(x)}+\ket{1\oplus f(x)}\big).
\]
Whether $f(x)=0$ or $f(x)=1$, the set $\{f(x),\,1\oplus f(x)\}$ equals $\{0,1\}$,
so the target factor is always $\tfrac{1}{\sqrt2}(\ket0+\ket1)=\ket{+}$ with a
$+$ sign. Thus $\Uf\ket{x}\ket{+}=\ket{x}\ket{+}$ and no phase depends on $f(x)$:
$\ket{+}$ is a $+1$ eigenstate of the target's bit-flip action. The state $\ket{-}$
is the $-1$ eigenstate, so flipping (which happens exactly when $f(x)=1$) multiplies
it by $-1$; that sign is what encodes $f$. Hence $\ket{-}$ is the useful choice.
\end{solution}

\begin{solution}[Exercise~\ref{l3:ex:p4}]
Here $f(00)=0,\ f(01)=1,\ f(10)=1,\ f(11)=0$: two zeros and two ones, so $f$ is
balanced. Starting from $\ket{00}$, $H^{\otimes2}$ gives
$\tfrac12(\ket{00}+\ket{01}+\ket{10}+\ket{11})$. Phase kickback multiplies each
$\ket{x}$ by $(-1)^{f(x)}$:
$\tfrac12(\ket{00}-\ket{01}-\ket{10}+\ket{11})$. Note this factors as
$\tfrac12(\ket0-\ket1)(\ket0-\ket1)=\ket{-}\ket{-}$. Applying $H^{\otimes2}$ gives
$H\ket{-}\otimes H\ket{-}=\ket1\ket1=\ket{11}$. The amplitude of $\ket{00}$ is $0$,
consistent with Theorem~\ref{l3:thm:dj}; indeed
$\alpha_{00}=\tfrac14\sum_x(-1)^{f(x)}=\tfrac14(1-1-1+1)=0$.
\end{solution}

\begin{solution}[Exercise~\ref{l3:ex:p5}]
With $g(x)=s\cdot x\oplus c$, phase kickback gives each branch the sign
$(-1)^{g(x)}=(-1)^{c}(-1)^{s\cdot x}$. The factor $(-1)^c$ is a \emph{global} phase
independent of $x$, so it can be pulled out of the sum:
\[
  \tfrac{1}{\sqrt N}\sum_x(-1)^{g(x)}\ket{x}
  =(-1)^{c}\,\tfrac{1}{\sqrt N}\sum_x(-1)^{s\cdot x}\ket{x}.
\]
Applying $H^{\otimes n}$ yields $(-1)^{c}\ket{s}$ by Theorem~\ref{l3:thm:bv}. A global
phase is physically undetectable and does not affect measurement probabilities, so
the outcome is $s$ with certainty. The constant $c$ is invisible to this algorithm;
recovering $c$ would require an additional query, e.g.\ evaluating $g(0^n)=c$.
\end{solution}

\begin{solution}[Exercise~\ref{l3:ex:p6}]
The inner product splits over bits, so the sum factors:
\[
  \frac{1}{2^n}\sum_{x\in\{0,1\}^n}(-1)^{x\cdot w}
  =\prod_{i=1}^{n}\Big(\tfrac12\sum_{x_i\in\{0,1\}}(-1)^{x_i w_i}\Big).
\]
For each factor, if $w_i=0$ the summand is $(-1)^0=1$ twice, giving
$\tfrac12(1+1)=1$; if $w_i=1$ the summand is $+1$ and $-1$, giving
$\tfrac12(1-1)=0$. The product is $1$ only when every $w_i=0$, i.e.\ $w=0^n$, and
$0$ otherwise. Hence the sum equals $\delta_{w,0^n}$.
\end{solution}

\begin{solution}[Exercise~\ref{l3:ex:p7}]
Deterministically, an adversary answering queries can keep all responses equal
(say all $0$) for the first $2^{n-1}$ distinct inputs; at that point $f$ could still
be constant (all outputs $0$) or balanced (the remaining $2^{n-1}$ inputs all $1$).
Only a $(2^{n-1}+1)$-th query can force the decision, and $2^{n-1}+1$ queries always
suffice: if any two of them disagree, $f$ is balanced, otherwise the majority
argument forces it to be constant. So the exact worst-case deterministic query
complexity is $2^{n-1}+1$. Randomly, sampling $k$ distinct inputs and answering
``balanced'' iff any two disagree: a constant $f$ is never misclassified, while a
balanced $f$ is missed only if all $k$ samples share a value, an event of
probability $2\cdot 2^{-k}$ (up to the small correction for sampling without
replacement). Taking $k=3$ already gives error below $1/3$, so a constant number of
queries suffices with bounded error---this is why the exponential separation is
special to the deterministic model, as noted in Remark~\ref{l3:rem:significance}.
\end{solution}

\chapter{Grover's Search Algorithm}

% The braket package ships \ketbra commented out; define it if the preamble has not.
\providecommand{\ketbra}[2]{\ket{#1}\!\bra{#2}}

Grover's algorithm solves the \emph{unstructured search} problem with a
quadratic speedup over any classical method. Given oracle access to a Boolean
function that flags a unique marked item among $N=2^{n}$ candidates, a classical
computer must, in the worst case, probe essentially all of them; Grover's
algorithm finds the marked item using only $O(\sqrt N)$ oracle queries. The
entire mechanism reduces to a strikingly simple geometric picture: the search
takes place inside a two-dimensional plane, and each iteration is a rotation
toward the solution. This chapter builds that picture from the ground up, proves
the key facts, computes the optimal number of iterations, and closes with the
generalization to amplitude amplification and the matching $\Omega(\sqrt N)$
lower bound that shows the algorithm is optimal.

\section{The Unstructured Search Problem}

We begin by fixing the problem and the model of computation. The word
\emph{unstructured} is essential: there is no order, index, or exploitable
regularity among the items, so the trick behind binary search is unavailable.

\begin{definition}[Unstructured search]
Let $N=2^{n}$ and let $f\colon\{0,1\}^{n}\to\{0,1\}$ be a Boolean function to
which we have only \emph{oracle} access. We are promised that there is a unique
marked item $x_0\in\{0,1\}^{n}$ with
\[
  f(x)=\begin{cases}1,& x=x_0,\\ 0,& x\neq x_0.\end{cases}
\]
The oracle is the unitary $\Uf$ acting on $n+1$ qubits by
\[
  \Uf\ket{x}\ket{y}=\ket{x}\ket{y\oplus f(x)},
  \qquad x\in\{0,1\}^{n},\ y\in\{0,1\}.
\]
The goal is to determine $x_0$ using as few queries to $\Uf$ as possible.
\end{definition}

Each application of $\Uf$ counts as one \emph{query}, and we measure the cost of
an algorithm by its query complexity. Because the function has no structure, the
oracle is the only way to learn anything about $x_0$: the string $x_0$ cannot be
read off directly.

\begin{proposition}[Classical lower bound]
Any classical (possibly randomized) algorithm that identifies $x_0$ with constant
success probability must make $\Omega(N)$ queries to $f$ in the worst case.
\end{proposition}

\begin{proof}[Sketch]
A classical algorithm learns nothing about $x_0$ from a query $f(x)$ unless
$x=x_0$. An adversary may answer $f(x)=0$ for every queried $x$ until only one
unqueried candidate remains, which is consistent with $x_0$ being that last item.
Thus to be certain, an algorithm must probe $N-1$ items in the worst case. A
randomized algorithm querying a uniformly random subset of size $k$ finds $x_0$
with probability $k/N$, so achieving constant success probability requires
$k=\Omega(N)$ queries.
\end{proof}

This abstract task captures many concrete situations: searching an unsorted
database, deciding satisfiability of a constraint problem (treat any satisfying
assignment as a marked item and let the oracle test assignments), or inverting a
function and finding collisions. In every case the common structure is the same:
\emph{package ``is this candidate a solution?'' into an oracle}, and the classical
exhaustive cost is $O(N)$. Grover's algorithm attacks all of them at once.

\subsection{The phase oracle}

It is convenient to convert the bit-flip oracle $\Uf$ into a \emph{phase} oracle.
Prepare the answer qubit in the state
$\ket{-}=\tfrac{1}{\sqrt2}\bigl(\ket{0}-\ket{1}\bigr)$. Then
\[
  \Uf\ket{x}\ket{-}
    =\ket{x}\,\tfrac{1}{\sqrt2}\bigl(\ket{0\oplus f(x)}-\ket{1\oplus f(x)}\bigr)
    =(-1)^{f(x)}\ket{x}\ket{-}.
\]
The answer qubit is unchanged and factors out; on the $n$-qubit register the
oracle acts as $\ket{x}\mapsto(-1)^{f(x)}\ket{x}$. This phase-kickback trick lets
us treat the oracle as an $n$-qubit operator that simply attaches a minus sign to
the marked basis state.

\section{The Two Reflections and the Grover Iterate}

Two operators drive the algorithm. Both are \emph{reflections}, and their product
is a rotation. We first define them, prove they are genuine reflections, and then
assemble the Grover iterate.

Throughout, let
\[
  \ket{\psi_0}=H^{\otimes n}\ket{0^{n}}=\frac{1}{\sqrt N}\sum_{x\in\{0,1\}^{n}}\ket{x}
\]
be the uniform superposition over all $N$ basis states. This is the algorithm's
starting state; every candidate carries amplitude $1/\sqrt N$.

\begin{definition}[Marking reflection and diffusion operator]
The \emph{marking reflection} (phase oracle) is
\[
  R_{x_0}=\id-2\ketbra{x_0}{x_0},
\]
which flips the sign of the amplitude on $\ket{x_0}$ and leaves the orthogonal
complement fixed. The \emph{diffusion operator} is
\[
  R_{\psi_0}=2\ketbra{\psi_0}{\psi_0}-\id,
\]
which fixes $\ket{\psi_0}$ and flips the sign of everything orthogonal to it.
It can be written equivalently as
\[
  R_{\psi_0}=H^{\otimes n}\bigl(2\ketbra{0^{n}}{0^{n}}-\id\bigr)H^{\otimes n},
\]
which is how it is implemented as a circuit.
\end{definition}

\begin{lemma}[Both operators are reflections]
\label{l4:lem:reflections}
Each of $R_{x_0}$ and $R_{\psi_0}$ is Hermitian, unitary, and an involution
($R^2=\id$); geometrically each is a reflection of the Hilbert space about a
hyperplane.
\end{lemma}

\begin{proof}
Let $\ket{u}$ be a unit vector and set $P=\ketbra{u}{u}$, so $P=P^{\dagger}$ and
$P^{2}=P$ (an orthogonal projector). Consider $R=\id-2P$. Then
\[
  R^{\dagger}=\id-2P^{\dagger}=\id-2P=R,
\]
so $R$ is Hermitian. Moreover
\[
  R^{2}=(\id-2P)^{2}=\id-4P+4P^{2}=\id-4P+4P=\id,
\]
so $R$ is an involution. A Hermitian involution is unitary, since
$R^{\dagger}R=R^{2}=\id$. Its action is $R\ket{u}=-\ket{u}$ and $R\ket{w}=\ket{w}$
for every $\ket{w}\perp\ket{u}$, which is exactly a reflection about the hyperplane
$\ket{u}^{\perp}$.

Applying this with $\ket{u}=\ket{x_0}$ gives $R_{x_0}=\id-2\ketbra{x_0}{x_0}$
directly. For the diffusion operator, note
$R_{\psi_0}=2\ketbra{\psi_0}{\psi_0}-\id=-\bigl(\id-2\ketbra{\psi_0}{\psi_0}\bigr)$;
by the computation above $\id-2\ketbra{\psi_0}{\psi_0}$ is a Hermitian unitary
involution, and negating it preserves all three properties. Explicitly
$R_{\psi_0}\ket{\psi_0}=\ket{\psi_0}$ and $R_{\psi_0}\ket{w}=-\ket{w}$ for
$\ket{w}\perp\ket{\psi_0}$, so $R_{\psi_0}$ is the reflection through the axis
$\ket{\psi_0}$.
The identity
$R_{\psi_0}=H^{\otimes n}(2\ketbra{0^{n}}{0^{n}}-\id)H^{\otimes n}$ follows because
$H^{\otimes n}$ is unitary and $H^{\otimes n}\ket{0^{n}}=\ket{\psi_0}$.
\end{proof}

Intuitively $R_{\psi_0}$ implements an ``inversion about the mean'': writing a
state as $\sum_x a_x\ket{x}$ with mean amplitude $\bar a=\tfrac1N\sum_x a_x$, one
checks that $R_{\psi_0}\sum_x a_x\ket{x}=\sum_x(2\bar a-a_x)\ket{x}$, reflecting
each amplitude about $\bar a$.

\begin{definition}[Grover iterate]
The \emph{Grover iterate} is the composition of the two reflections,
\[
  G=R_{\psi_0}R_{x_0}
   =\bigl(2\ketbra{\psi_0}{\psi_0}-\id\bigr)\bigl(\id-2\ketbra{x_0}{x_0}\bigr).
\]
One iteration first marks the solution (via $R_{x_0}$) and then reflects about the
mean (via $R_{\psi_0}$).
\end{definition}

\section{The Invariant Plane and the Rotation Picture}

The power of the algorithm comes from the fact that $G$ never leaves a fixed
two-dimensional plane, inside which it acts as an ordinary rotation.

Define the normalized state orthogonal to $\ket{x_0}$ that lives in the ``bad''
subspace:
\[
  \ket{\psi_0^{\perp}}
    =\frac{1}{\sqrt{N-1}}\sum_{x\neq x_0}\ket{x}.
\]
Then $\braket{x_0|\psi_0^{\perp}}=0$, and the uniform state decomposes as
\begin{equation}
  \ket{\psi_0}
    =\frac{1}{\sqrt N}\ket{x_0}+\frac{\sqrt{N-1}}{\sqrt N}\ket{\psi_0^{\perp}}
    =\sin\frac{\theta}{2}\,\ket{x_0}+\cos\frac{\theta}{2}\,\ket{\psi_0^{\perp}},
  \label{l4:eq:initdecomp}
\end{equation}
where the angle $\theta\in(0,\pi)$ is defined by
\begin{equation}
  \sin\frac{\theta}{2}=\frac{1}{\sqrt N},
  \qquad
  \theta=2\arcsin\frac{1}{\sqrt N}\approx\frac{2}{\sqrt N}\quad(N\gg1).
  \label{l4:eq:theta}
\end{equation}
So $\ket{\psi_0}$ makes a small angle $\theta/2$ with the axis
$\ket{\psi_0^{\perp}}$.

\begin{theorem}[$G$ is a rotation in the plane $\mathcal{P}$]
\label{l4:thm:rotation}
Let $\mathcal{P}=\operatorname{span}\{\ket{x_0},\ket{\psi_0^{\perp}}\}$. Then $G$
maps $\mathcal{P}$ into itself, and restricted to $\mathcal{P}$ (in the ordered
orthonormal basis $\{\ket{\psi_0^{\perp}},\ket{x_0}\}$) it is the rotation by the
angle $\theta$ of \eqref{l4:eq:theta}:
\[
  G\big|_{\mathcal{P}}
    =\begin{pmatrix}\cos\theta & -\sin\theta\\[2pt]\sin\theta & \cos\theta\end{pmatrix}.
\]
\end{theorem}

\begin{proof}
First note that both reflections preserve $\mathcal{P}$. The vector $\ket{x_0}$
lies in $\mathcal{P}$ and $\ket{\psi_0}\in\mathcal{P}$ by \eqref{l4:eq:initdecomp},
so $\ket{\psi_0^{\perp}}\in\mathcal{P}$ as well. The operator $R_{x_0}$ negates
$\ket{x_0}$ and fixes every vector orthogonal to $\ket{x_0}$; in particular it
fixes $\ket{\psi_0^{\perp}}$, so it maps $\mathcal{P}$ to itself. The operator
$R_{\psi_0}$ fixes $\ket{\psi_0}\in\mathcal P$ and negates its orthogonal
complement; the in-plane vector orthogonal to $\ket{\psi_0}$ inside $\mathcal P$
stays in $\mathcal P$, so $R_{\psi_0}$ also preserves $\mathcal{P}$. Hence
$G=R_{\psi_0}R_{x_0}$ preserves $\mathcal{P}$.

We now use the plane-geometry fact that \emph{the product of two reflections is a
rotation by twice the angle between their axes}. Inside $\mathcal P$, the mirror
line of $R_{x_0}$ is the axis $\ket{\psi_0^{\perp}}$ (the fixed direction of
$R_{x_0}$ within $\mathcal P$), and the mirror line of $R_{\psi_0}$ is the axis
$\ket{\psi_0}$. By \eqref{l4:eq:initdecomp} the angle between $\ket{\psi_0^{\perp}}$
and $\ket{\psi_0}$ is $\theta/2$, so $G$ is a rotation by $2\cdot(\theta/2)=\theta$.

For a self-contained verification, compute the matrix directly. Using
$\sin\tfrac{\theta}{2}=1/\sqrt N$, $\cos\tfrac{\theta}{2}=\sqrt{(N-1)/N}$ and
\eqref{l4:eq:initdecomp}, apply $R_{x_0}=\id-2\ketbra{x_0}{x_0}$:
\[
  R_{x_0}\ket{\psi_0^{\perp}}=\ket{\psi_0^{\perp}},
  \qquad
  R_{x_0}\ket{x_0}=-\ket{x_0}.
\]
Thus in the basis $\{\ket{\psi_0^{\perp}},\ket{x_0}\}$,
$R_{x_0}=\begin{psmallmatrix}1&0\\0&-1\end{psmallmatrix}$, a reflection about the
$\ket{\psi_0^{\perp}}$ axis. Next $R_{\psi_0}=2\ketbra{\psi_0}{\psi_0}-\id$ is the
reflection about the axis $\ket{\psi_0}$, which sits at angle $\theta/2$ from
$\ket{\psi_0^{\perp}}$; a reflection about a line at angle $\phi$ from the first
basis vector has matrix
$\begin{psmallmatrix}\cos2\phi&\sin2\phi\\\sin2\phi&-\cos2\phi\end{psmallmatrix}$,
so with $\phi=\theta/2$,
\[
  R_{\psi_0}=\begin{pmatrix}\cos\theta&\sin\theta\\ \sin\theta&-\cos\theta\end{pmatrix}.
\]
Multiplying,
\[
  G=R_{\psi_0}R_{x_0}
   =\begin{pmatrix}\cos\theta&\sin\theta\\ \sin\theta&-\cos\theta\end{pmatrix}
    \begin{pmatrix}1&0\\0&-1\end{pmatrix}
   =\begin{pmatrix}\cos\theta&-\sin\theta\\ \sin\theta&\cos\theta\end{pmatrix},
\]
which is exactly the rotation by $\theta$.
\end{proof}

\begin{theorem}[State after $k$ iterations]
\label{l4:thm:afterk}
With $\theta$ as in \eqref{l4:eq:theta},
\[
  G^{k}\ket{\psi_0}
    =\sin\!\frac{(2k+1)\theta}{2}\,\ket{x_0}
     +\cos\!\frac{(2k+1)\theta}{2}\,\ket{\psi_0^{\perp}}.
\]
Consequently the probability of measuring the marked item after $k$ iterations is
$p_k=\sin^{2}\!\bigl((2k+1)\theta/2\bigr)$.
\end{theorem}

\begin{proof}
By Theorem~\ref{l4:thm:rotation}, $G$ restricted to $\mathcal P$ is rotation by
$\theta$, so $G^{k}$ is rotation by $k\theta$. The initial state
\eqref{l4:eq:initdecomp} sits at angle $\theta/2$ measured from the
$\ket{\psi_0^{\perp}}$ axis (its $\ket{x_0}$-component is
$\sin(\theta/2)$). Rotating by $k\theta$ increases this angle to
$\theta/2+k\theta=(2k+1)\theta/2$. Hence the $\ket{x_0}$-component becomes
$\sin\bigl((2k+1)\theta/2\bigr)$ and the $\ket{\psi_0^{\perp}}$-component becomes
$\cos\bigl((2k+1)\theta/2\bigr)$, as claimed. Measuring in the computational basis
yields $x_0$ with probability equal to the squared $\ket{x_0}$-amplitude, namely
$p_k=\sin^{2}\bigl((2k+1)\theta/2\bigr)$.
\end{proof}

\section{Optimal Iteration Count and Query Complexity}

Theorem~\ref{l4:thm:afterk} shows the amplitude on $\ket{x_0}$ is maximized when the
angle $(2k+1)\theta/2$ is as close as possible to $\pi/2$, i.e.\ when the state is
(nearly) aligned with $\ket{x_0}$.

\begin{theorem}[Optimal number of iterations]
\label{l4:thm:kstar}
Choose $k$ so that $(2k+1)\theta/2$ is closest to $\pi/2$; the optimal choice is
\[
  k^{\ast}=\left\lfloor\frac{\pi}{4}\sqrt N\right\rfloor
  \qquad\bigl(\text{rounding } \tfrac{\pi}{2\theta}-\tfrac12\bigr).
\]
With this choice the success probability satisfies $p_{k^{\ast}}\ge 1-1/N$, so a
single run succeeds with probability $\approx1$. The algorithm therefore uses
$k^{\ast}=O(\sqrt N)$ oracle queries.
\end{theorem}

\begin{proof}
We want $(2k+1)\theta/2\approx\pi/2$, i.e.\ $2k+1\approx\pi/\theta$, i.e.\
$k\approx\pi/(2\theta)-1/2$. Using the approximation
$\theta\approx2/\sqrt N$ from \eqref{l4:eq:theta} valid for large $N$,
\[
  k\approx\frac{\pi}{2\theta}-\frac12\approx\frac{\pi}{2}\cdot\frac{\sqrt N}{2}-\frac12
     =\frac{\pi}{4}\sqrt N-\frac12,
\]
and rounding to the nearest nonnegative integer gives
$k^{\ast}=\lfloor(\pi/4)\sqrt N\rfloor$. Because $k$ is an integer we cannot in
general hit the angle $\pi/2$ exactly, but the residual angular error is at most
$\theta/2$ (half a rotation step): choosing $k^{\ast}$ so that
$(2k^{\ast}+1)\theta/2$ lies within $\theta/2$ of $\pi/2$ gives
\[
  p_{k^{\ast}}=\sin^{2}\!\frac{(2k^{\ast}+1)\theta}{2}
    =\cos^{2}\!\Bigl(\frac{\pi}{2}-\frac{(2k^{\ast}+1)\theta}{2}\Bigr)
    \ge\cos^{2}\!\frac{\theta}{2}=1-\sin^{2}\frac{\theta}{2}=1-\frac1N.
\]
Each iterate $G=R_{\psi_0}R_{x_0}$ makes exactly one oracle call (inside
$R_{x_0}$), so the total query cost is $k^{\ast}=O(\sqrt N)$.
\end{proof}

This is the headline result: $O(\sqrt N)$ queries versus the classical $O(N)$, a
quadratic speedup. The gain comes from operating on \emph{amplitudes} rather than
probabilities: $G$ rotates the amplitude linearly in $k$, and probability is its
square, so the linear growth of the angle produces the square-root scaling in $k$.

\begin{remark}[Do not over-rotate]
More iterations is \emph{not} always better. Since $p_k=\sin^{2}((2k+1)\theta/2)$
oscillates, once $(2k+1)\theta/2$ passes $\pi/2$ the success probability decreases
again --- the state ``rotates past'' $\ket{x_0}$. One must therefore stop at
$k^{\ast}$; running the loop too long is counterproductive. If $N$ is unknown one
can still succeed with $O(\sqrt N)$ expected queries by trying geometrically
increasing guesses for the number of iterations and choosing the number at random
within each guess.
\end{remark}

\section{Multiple Solutions and Amplitude Amplification}

Suppose now that there are $M\ge1$ marked items, $S=\{x:f(x)=1\}$ with $\abs{S}=M$.
Everything goes through with the ``good'' subspace replaced by the span of the
marked basis states.

Let
\[
  \ket{\mathrm{good}}=\frac{1}{\sqrt M}\sum_{x\in S}\ket{x},
  \qquad
  \ket{\mathrm{bad}}=\frac{1}{\sqrt{N-M}}\sum_{x\notin S}\ket{x},
\]
so that $\ket{\psi_0}=\sqrt{M/N}\,\ket{\mathrm{good}}+\sqrt{(N-M)/N}\,\ket{\mathrm{bad}}$.
Define $R_S=\id-2\sum_{x\in S}\ketbra{x}{x}$ (the phase oracle now flips all marked
items) and keep the same diffusion $R_{\psi_0}$. The identical argument shows
$G=R_{\psi_0}R_S$ is a rotation by $\theta$ in
$\operatorname{span}\{\ket{\mathrm{good}},\ket{\mathrm{bad}}\}$ with
\[
  \sin\frac{\theta}{2}=\sqrt{\frac{M}{N}},
  \qquad\text{equivalently}\qquad
  \sin\theta=2\sin\frac{\theta}{2}\cos\frac{\theta}{2}
            =2\sqrt{\frac{M}{N}}\sqrt{\frac{N-M}{N}}\approx 2\sqrt{\frac{M}{N}}.
\]
The optimal iteration count becomes
\[
  k^{\ast}=\left\lfloor\frac{\pi}{4}\sqrt{\frac{N}{M}}\right\rfloor
  =O\!\left(\sqrt{\frac{N}{M}}\right),
\]
so more solutions make the search faster.

\begin{remark}[Amplitude amplification]
Grover's construction generalizes far beyond search. Suppose $\mathcal A$ is any
(quantum) procedure preparing a state
$\mathcal A\ket{0}=\sqrt{p_0}\,\ket{\psi_{\mathrm{good}}}+\sqrt{1-p_0}\,\ket{\psi_{\mathrm{bad}}}$,
where a projector $P$ recognizes the good part and $p_0$ is the initial success
probability. Replacing $R_{x_0}$ by the reflection $\id-2P$ and $R_{\psi_0}$ by
$\mathcal A(2\ketbra{0}{0}-\id)\mathcal A^{\dagger}$ yields an iterate that rotates
by an angle $\theta$ with $\sin^{2}(\theta/2)=p_0$. After $O(1/\sqrt{p_0})$
iterations the good state is reached with probability $\Omega(1)$ --- a quadratic
improvement over the $O(1/p_0)$ repetitions that classical
prepare-and-check would require. Grover search is the special case
$p_0=M/N$.
\end{remark}

\begin{example}[The exact case $N=4$, $M=1$]
Here $\sin(\theta/2)=1/\sqrt4=1/2$, so $\theta/2=\pi/6$ and $\theta=\pi/3$. After a
single iteration ($k=1$) the angle is
\[
  \frac{(2\cdot1+1)\theta}{2}=\frac{3\theta}{2}=\frac{3}{2}\cdot\frac{\pi}{3}=\frac{\pi}{2},
\]
so by Theorem~\ref{l4:thm:afterk} the state is
$G\ket{\psi_0}=\sin(\pi/2)\ket{x_0}+\cos(\pi/2)\ket{\psi_0^{\perp}}=\ket{x_0}$.
Thus for $N=4$ a \emph{single} Grover iteration finds the marked item with
certainty --- a rare instance where the algorithm is exactly deterministic.
\end{example}

\begin{example}[A multi-solution estimate]
Suppose a SAT instance on $n=30$ variables has $M=2^{10}$ satisfying assignments
among $N=2^{30}$ candidates. Classical exhaustive search needs on the order of
$N/M=2^{20}\approx10^{6}$ trials to hit a solution. Grover instead needs about
\[
  \frac{\pi}{4}\sqrt{\frac{N}{M}}
    =\frac{\pi}{4}\sqrt{2^{20}}
    =\frac{\pi}{4}\cdot2^{10}\approx804
\]
oracle queries --- the square-root improvement in action.
\end{example}

\section{Optimality: the \texorpdfstring{$\Omega(\sqrt N)$}{Omega(sqrt N)} Lower Bound}

Grover's algorithm is not just good; it is asymptotically the best possible.

\begin{theorem}[Query lower bound]
Any quantum algorithm that finds the marked item $x_0$ (or even merely decides
whether a marked item exists) with constant success probability must make
$\Omega(\sqrt N)$ oracle queries.
\end{theorem}

\begin{proof}[Idea]
Run the algorithm both on the ``empty'' oracle $f\equiv0$ and on the oracle marking
a single item $x_0$. Let $\ket{\phi_k}$ be the state after $k$ queries in the
empty case and $\ket{\phi_k^{x_0}}$ the state in the marked case, and define the
total deviation
\[
  D_k=\sum_{x_0}\norm{\,\ket{\phi_k^{x_0}}-\ket{\phi_k}\,}^{2}.
\]
A hybrid argument shows each additional query can increase $D_k$ by only a bounded
amount: $D_k\le 4k^{2}$. On the other hand, to reliably distinguish which item is
marked, the marked-case states for different $x_0$ must be nearly orthogonal to the
empty-case state, which forces $D_k=\Omega(N)$. Combining
$\Omega(N)\le D_k\le 4k^{2}$ gives $k=\Omega(\sqrt N)$.
\end{proof}

\begin{remark}[Grover is optimal]
The upper bound of Theorem~\ref{l4:thm:kstar} and the lower bound above match up to
constant factors: both are $\Theta(\sqrt N)$. Hence Grover's algorithm is optimal
for unstructured search --- no quantum algorithm can do asymptotically better. The
quadratic speedup is the best achievable when the only access to the problem is
through an unstructured oracle; a super-polynomial speedup would require exploitable
structure, as in Shor's algorithm or the phase-estimation-based methods.
\end{remark}

\section{Problems}

\begin{exercise}
\label{l4:ex:invmean}
Show that the diffusion operator implements ``inversion about the mean'': for any
state $\sum_x a_x\ket{x}$ with mean amplitude $\bar a=\tfrac1N\sum_x a_x$,
\[
  R_{\psi_0}\sum_x a_x\ket{x}=\sum_x(2\bar a-a_x)\ket{x}.
\]
\end{exercise}

\begin{exercise}
\label{l4:ex:phaseoracle}
Verify the phase-kickback identity
\[
  \Uf\ket{x}\ket{-}=(-1)^{f(x)}\ket{x}\ket{-}
\]
starting from $\Uf\ket{x}\ket{y}=\ket{x}\ket{y\oplus f(x)}$, and explain why this
lets us regard the oracle as the $n$-qubit operator $R_{x_0}$.
\end{exercise}

\begin{exercise}
\label{l4:ex:n2}
Take $N=2$ (one qubit, $n=1$) with the single marked item $x_0$. Compute $\theta$
from \eqref{l4:eq:theta}, apply one Grover iterate to $\ket{\psi_0}$, and interpret the
result. Why is the case $N=2$ slightly special?
\end{exercise}

\begin{exercise}
\label{l4:ex:success}
For general $N$, using $k^{\ast}=\lfloor(\pi/4)\sqrt N\rfloor$, show that the
success probability satisfies $p_{k^{\ast}}\ge1-1/N$. (You may use that the
residual angle from $\pi/2$ is at most $\theta/2$.)
\end{exercise}

\begin{exercise}
\label{l4:ex:overrotate}
Consider $N=2^{20}$. Approximately how many iterations $k^{\ast}$ maximize the
success probability? What happens qualitatively to the success probability if one
instead runs $2k^{\ast}$ iterations? Explain using $p_k=\sin^2((2k+1)\theta/2)$.
\end{exercise}

\begin{exercise}
\label{l4:ex:multi}
There are $M=4$ marked items among $N=1024$. Find $\sin(\theta/2)$, the optimal
iteration count $k^{\ast}$, and compare with the single-solution count for the same
$N$.
\end{exercise}

\begin{exercise}
\label{l4:ex:reflprod}
Prove the plane-geometry fact used in Theorem~\ref{l4:thm:rotation}: in a real
two-dimensional plane, the composition of a reflection about a line $\ell_1$ and a
reflection about a line $\ell_2$ (with $\ell_2$ at angle $\alpha$ from $\ell_1$) is
a rotation by $2\alpha$.
\end{exercise}

\subsection*{Solutions}

\begin{solution}[Exercise~\ref{l4:ex:invmean}]
Write $R_{\psi_0}=2\ketbra{\psi_0}{\psi_0}-\id$ and let $\ket{\phi}=\sum_x a_x\ket{x}$.
Since $\ket{\psi_0}=\tfrac1{\sqrt N}\sum_x\ket{x}$, we have
$\braket{\psi_0|\phi}=\tfrac1{\sqrt N}\sum_x a_x=\sqrt N\,\bar a$. Therefore
\[
  2\ketbra{\psi_0}{\psi_0}\phi
    =2\sqrt N\,\bar a\,\ket{\psi_0}
    =2\sqrt N\,\bar a\cdot\frac1{\sqrt N}\sum_x\ket{x}
    =\sum_x 2\bar a\,\ket{x}.
\]
Subtracting $\id\ket{\phi}=\sum_x a_x\ket{x}$ gives
$R_{\psi_0}\ket{\phi}=\sum_x(2\bar a-a_x)\ket{x}$, i.e.\ each amplitude is reflected
about the mean $\bar a$. This is why $R_{\psi_0}$ is often called ``inversion about
the mean.''
\end{solution}

\begin{solution}[Exercise~\ref{l4:ex:phaseoracle}]
With $\ket{-}=\tfrac1{\sqrt2}(\ket0-\ket1)$,
\[
  \Uf\ket{x}\ket{-}
   =\tfrac1{\sqrt2}\bigl(\ket{x}\ket{0\oplus f(x)}-\ket{x}\ket{1\oplus f(x)}\bigr).
\]
If $f(x)=0$ this is $\tfrac1{\sqrt2}(\ket{x}\ket0-\ket{x}\ket1)=\ket{x}\ket{-}$. If
$f(x)=1$ the two kets swap, giving
$\tfrac1{\sqrt2}(\ket{x}\ket1-\ket{x}\ket0)=-\ket{x}\ket{-}$. Both cases combine into
$(-1)^{f(x)}\ket{x}\ket{-}$. The answer qubit stays in $\ket{-}$ and factors out, so
on the $n$-qubit register the oracle acts as $\ket{x}\mapsto(-1)^{f(x)}\ket{x}$,
which negates exactly $\ket{x_0}$ --- that is $R_{x_0}=\id-2\ketbra{x_0}{x_0}$.
\end{solution}

\begin{solution}[Exercise~\ref{l4:ex:n2}]
For $N=2$, $\sin(\theta/2)=1/\sqrt2$, so $\theta/2=\pi/4$ and $\theta=\pi/2$. The
initial state is $\ket{\psi_0}=\tfrac1{\sqrt2}(\ket{x_0}+\ket{\psi_0^{\perp}})$, at
angle $\theta/2=\pi/4$ from the $\ket{\psi_0^{\perp}}$ axis. One iterate rotates by
$\theta=\pi/2$, bringing the angle to $\pi/4+\pi/2=3\pi/4$; equivalently
$(2\cdot1+1)\theta/2=3\pi/4$, and $p_1=\sin^2(3\pi/4)=1/2$. So a single iterate does
\emph{not} give certainty here, and indeed
$\lfloor(\pi/4)\sqrt2\rfloor=\lfloor1.11\rfloor=1$ while the best achievable
probability for $N=2$ is $1/2$, no better than a random guess. The case is special
because with only two items the initial angle $\pi/4$ is already large and the
rotation overshoots; unstructured search offers no advantage when $N=2$.
\end{solution}

\begin{solution}[Exercise~\ref{l4:ex:success}]
Pick $k^{\ast}$ so that the angle $(2k^{\ast}+1)\theta/2$ is within $\theta/2$ of
$\pi/2$; this is possible because consecutive angles differ by $\theta$. Writing the
residual $\delta=\pi/2-(2k^{\ast}+1)\theta/2$ with $\abs{\delta}\le\theta/2$,
\[
  p_{k^{\ast}}=\sin^2\!\Bigl(\tfrac\pi2-\delta\Bigr)=\cos^2\delta\ge\cos^2\tfrac\theta2
   =1-\sin^2\tfrac\theta2=1-\tfrac1N,
\]
using $\sin(\theta/2)=1/\sqrt N$. Hence the success probability is at least
$1-1/N\to1$, so one run essentially always succeeds; if it fails one verifies with a
single extra oracle call and repeats.
\end{solution}

\begin{solution}[Exercise~\ref{l4:ex:overrotate}]
For $N=2^{20}=1{,}048{,}576$, $\sqrt N=1024$, so
$k^{\ast}=\lfloor(\pi/4)\cdot1024\rfloor=\lfloor 804.2\rfloor=804$. At $k^{\ast}$ the
angle $(2k^{\ast}+1)\theta/2\approx\pi/2$ and $p_{k^{\ast}}\approx1$. Running instead
$2k^{\ast}\approx1608$ iterations roughly doubles the angle to about
$(2\cdot2k^{\ast}+1)\theta/2\approx 3\pi/2$; since $p_k$ oscillates as $\sin^2$, the
state has rotated well past $\ket{x_0}$ and the success probability has fallen back
toward small values. The lesson: the probability oscillates in $k$, so over-rotating
destroys the amplitude gained --- stop at $k^{\ast}$.
\end{solution}

\begin{solution}[Exercise~\ref{l4:ex:multi}]
With $M=4$, $N=1024$: $\sin(\theta/2)=\sqrt{M/N}=\sqrt{4/1024}=\sqrt{1/256}=1/16$.
The optimal count is
$k^{\ast}=\lfloor(\pi/4)\sqrt{N/M}\rfloor=\lfloor(\pi/4)\sqrt{256}\rfloor
=\lfloor(\pi/4)\cdot16\rfloor=\lfloor12.57\rfloor=12$. For a single solution
($M=1$) with the same $N$ one would need
$\lfloor(\pi/4)\sqrt{1024}\rfloor=\lfloor(\pi/4)\cdot32\rfloor=\lfloor25.13\rfloor=25$
iterations. So having $M=4$ solutions roughly halves the query count (from $25$ to
$12$), consistent with the $O(\sqrt{N/M})$ scaling: a factor of $2$ speedup for a
factor of $4$ more solutions.
\end{solution}

\begin{solution}[Exercise~\ref{l4:ex:reflprod}]
Place coordinates so that $\ell_1$ is the $x$-axis and $\ell_2$ makes angle $\alpha$
with it. Reflection about a line at angle $\beta$ has matrix
$M(\beta)=\begin{psmallmatrix}\cos2\beta&\sin2\beta\\\sin2\beta&-\cos2\beta\end{psmallmatrix}$.
Thus reflection about $\ell_1$ is
$M(0)=\begin{psmallmatrix}1&0\\0&-1\end{psmallmatrix}$
and about $\ell_2$ is $M(\alpha)$. Their product is
\[
  M(\alpha)M(0)
   =\begin{pmatrix}\cos2\alpha&\sin2\alpha\\ \sin2\alpha&-\cos2\alpha\end{pmatrix}
    \begin{pmatrix}1&0\\0&-1\end{pmatrix}
   =\begin{pmatrix}\cos2\alpha&-\sin2\alpha\\ \sin2\alpha&\cos2\alpha\end{pmatrix},
\]
which is the rotation matrix by angle $2\alpha$. Hence composing the two reflections
gives a rotation by twice the angle between the mirror lines. In
Theorem~\ref{l4:thm:rotation} the two mirror axes are $\ket{\psi_0^{\perp}}$ and
$\ket{\psi_0}$, separated by $\theta/2$, giving rotation by $\theta$.
\end{solution}

\chapter{The Quantum Fourier Transform}

The Quantum Fourier Transform (QFT) is one of the most important subroutines in
quantum computing. It is the quantum analogue of the classical discrete Fourier
transform (DFT), and it sits at the heart of quantum phase estimation (QPE) and
of Shor's factoring algorithm. In this chapter we define the QFT, prove that it
is unitary, derive its product (tensor) representation, use that representation
to build an efficient circuit from Hadamard and controlled-phase gates, count
the gates to obtain the $O(n^2)$ bound, and discuss the inverse QFT together
with its principal applications.

Throughout we work with $n$ qubits, we write $N = 2^n$ for the dimension of the
state space $(\C^2)^{\otimes n}\cong\C^N$, and we abbreviate the primitive
$N$-th root of unity by
\[
  \omega \defeq e^{i 2\pi / N}, \qquad \omega^N = 1 .
\]
Computational basis states are labelled by integers $j\in\{0,1,\dots,N-1\}$, and
we freely identify the integer $j$ with its $n$-bit binary string
$j_1 j_2 \cdots j_n$, where $j_1$ is the most significant bit, so that
\[
  j = j_1 2^{n-1} + j_2 2^{n-2} + \cdots + j_n 2^{0}
    = \sum_{l=1}^{n} j_l\, 2^{\,n-l},
  \qquad \ket{j} = \ket{j_1}\otimes\cdots\otimes\ket{j_n}.
\]

\section{Definition and Motivation}

\begin{definition}[Quantum Fourier Transform]\label{l5:def:qft}
Let $N = 2^n$ and $\omega = e^{i 2\pi / N}$. The \emph{Quantum Fourier
Transform} is the linear operator $\QFT$ on $\C^N$ that acts on computational
basis states by
\[
  \QFT\ket{j} \;=\; \frac{1}{\sqrt N}\sum_{k=0}^{N-1} e^{i 2\pi j k / N}\ket{k}
             \;=\; \frac{1}{\sqrt N}\sum_{k=0}^{N-1} \omega^{jk}\,\ket{k},
  \qquad j = 0,1,\dots,N-1 .
\]
Equivalently, in the computational basis $\QFT$ is the $N\times N$ matrix with
entries
\[
  \bigl(\QFT\bigr)_{kj} \;=\; \frac{1}{\sqrt N}\,\omega^{jk}
  \qquad (0 \le j,k \le N-1),
\]
a (normalized) Vandermonde matrix in the powers of $\omega$.
\end{definition}

The action of $\QFT$ on a general state follows by linearity. If
$\ket{\psi} = \sum_{j=0}^{N-1} x_j \ket{j}$ then
\[
  \QFT\ket{\psi}
    = \sum_{k=0}^{N-1} y_k \ket{k},
  \qquad
  y_k = \frac{1}{\sqrt N}\sum_{j=0}^{N-1} \omega^{jk} x_j ,
\]
so the QFT transforms the vector of amplitudes $(x_0,\dots,x_{N-1})$ into
$(y_0,\dots,y_{N-1})$ exactly as the classical DFT would. The essential point is
that the input is a single computational basis state $\ket{j}$ while the output
is an \emph{equal-magnitude} superposition of all $N$ basis states; the integer
$j$ is encoded entirely in the relative phases $\omega^{jk}$.

\begin{example}[The all-zeros input]\label{l5:ex:zero}
For $j=0$ every phase is $\omega^{0}=1$, so
\[
  \QFT\ket{0} = \frac{1}{\sqrt N}\sum_{k=0}^{N-1}\ket{k}
              = H^{\otimes n}\ket{0}^{\otimes n},
\]
the uniform superposition. Thus on the input $\ket{0\cdots 0}$ the QFT is
nothing but a layer of Hadamard gates.
\end{example}

\begin{remark}[Same kernel as the classical DFT, but no free lunch]\label{l5:rem:dft}
The kernel $\tfrac{1}{\sqrt N}\,e^{i2\pi jk/N}$ appearing in
Definition~\ref{l5:def:qft} is \emph{identical} to that of the (symmetrically
normalized) classical DFT; the QFT simply applies this transform to the $2^n$
amplitudes of an $n$-qubit state \emph{in superposition}. It is tempting to
conclude that the QFT gives an exponential speedup for Fourier analysis, but
this is not so: after the transform we hold the coefficients $y_k$ only as
quantum amplitudes, and a measurement returns a \emph{single} basis state
$\ket{k}$ with probability $\abs{y_k}^2$. We cannot read out all $N$ amplitudes;
in fact estimating them to good precision would require exponentially many
repetitions. The power of the QFT is therefore not standalone but
\emph{structural}: embedded inside a larger algorithm (phase estimation,
period finding) it converts information hidden in the phases into information in
the measurement statistics, so that a small number of measurements suffices. We
also note the resource comparison developed below: the QFT circuit uses
$O(n^2) = O((\log N)^2)$ elementary gates, versus $O(N\log N)$ arithmetic
operations for the classical FFT and $O(N^2)$ for the naive DFT.
\end{remark}

\section{Unitarity}

Because the QFT is to be realized as a quantum circuit, it must be unitary. We
prove this directly from the definition. The engine of the proof is the
following elementary but crucial identity.

\begin{lemma}[Geometric-series / Kronecker-delta identity]\label{l5:lem:geom}
Let $N\ge 1$ and let $m\in\Z$. Then
\[
  \sum_{k=0}^{N-1} e^{i 2\pi m k / N}
  = \sum_{k=0}^{N-1}\bigl(e^{i2\pi m/N}\bigr)^{k}
  = N\,[\,m\equiv 0 \bmod N\,],
\]
where the bracket equals $1$ if $m\equiv 0 \pmod N$ and $0$ otherwise. In
particular, for $m = j - j'$ with $j,j' \in \{0,\dots,N-1\}$ we have
$m\equiv 0 \pmod N$ iff $j = j'$, so
\[
  \sum_{k=0}^{N-1} e^{i 2\pi (j-j') k / N} = N\,\delta_{j,j'} .
\]
\end{lemma}

\begin{proof}
Write $r = e^{i 2\pi m/N}$, so the sum is the geometric series
$\sum_{k=0}^{N-1} r^{k}$. If $m\equiv 0\pmod N$ then $r=1$ and every term is
$1$, giving the value $N$. Otherwise $r\ne 1$, and the closed form of the
geometric series applies:
\[
  \sum_{k=0}^{N-1} r^{k} = \frac{r^{N}-1}{r-1}
  = \frac{e^{i 2\pi m}-1}{r-1} = \frac{1-1}{r-1} = 0,
\]
since $r^{N} = e^{i 2\pi m} = 1$ for integer $m$ while the denominator is
nonzero. For the final statement, if $\abs{j-j'} < N$ then $j-j'\equiv 0
\pmod N$ forces $j-j'=0$.
\end{proof}

\begin{theorem}[The QFT is unitary]\label{l5:thm:unitary}
The operator of Definition~\ref{l5:def:qft} satisfies $\QFT^{\dagger}\QFT = \id$,
so $\QFT$ is unitary.
\end{theorem}

\begin{proof}
From Definition~\ref{l5:def:qft}, $\bigl(\QFT\bigr)_{kj} = N^{-1/2}\,\omega^{jk}$,
so the adjoint has entries
$\bigl(\QFT^{\dagger}\bigr)_{jk} = \overline{\bigl(\QFT\bigr)_{kj}}
= N^{-1/2}\,\omega^{-jk}$. We compute the matrix element
$\braket{j'|\QFT^{\dagger}\QFT|j}$, i.e.\ the $(j',j)$ entry of
$\QFT^{\dagger}\QFT$, by inserting a resolution of the identity over the
basis $\{\ket{k}\}$:
\[
  \braket{j'|\QFT^{\dagger}\QFT|j}
    = \sum_{k=0}^{N-1} \overline{(\QFT)_{kj'}}\,(\QFT)_{kj}
    = \sum_{k=0}^{N-1}
        \Bigl(\tfrac{1}{\sqrt N}\,\omega^{-j'k}\Bigr)
        \Bigl(\tfrac{1}{\sqrt N}\,\omega^{jk}\Bigr)
    = \frac{1}{N}\sum_{k=0}^{N-1}\omega^{(j-j')k}.
\]
Since $\omega^{(j-j')k} = e^{i2\pi (j-j')k/N}$, Lemma~\ref{l5:lem:geom} gives
$\sum_{k=0}^{N-1}\omega^{(j-j')k} = N\,\delta_{j,j'}$, and therefore
\[
  \braket{j'|\QFT^{\dagger}\QFT|j}
    = \frac{1}{N}\,N\,\delta_{j,j'} = \delta_{j,j'} .
\]
As $j,j'$ range over the basis this says $\QFT^{\dagger}\QFT = \id$. Equivalently,
the columns $\QFT\ket{j}$ are orthonormal. Hence $\QFT$ is unitary.
\end{proof}

Since $\QFT$ is unitary, $\QFT^{-1} = \QFT^{\dagger}$, and we record the inverse
explicitly.

\begin{definition}[Inverse QFT]\label{l5:def:iqft}
The \emph{inverse Quantum Fourier Transform} is $\QFT^{-1} = \QFT^{\dagger}$,
acting on basis states by
\[
  \QFT^{\dagger}\ket{j}
   = \frac{1}{\sqrt N}\sum_{k=0}^{N-1} e^{-i 2\pi j k / N}\ket{k}
   = \frac{1}{\sqrt N}\sum_{k=0}^{N-1} \omega^{-jk}\,\ket{k}.
\]
It differs from $\QFT$ only by conjugating every phase (replacing $\omega$ by
$\overline{\omega}=\omega^{-1}$). Consequently the inverse circuit is obtained
from the QFT circuit by reversing the order of the gates and negating every
phase angle.
\end{definition}

\section{The Product (Tensor) Representation}

To build a circuit we need the QFT output expressed qubit-by-qubit. The bridge
is the \emph{product form}, which shows that $\QFT\ket{j}$ is an unentangled
product state. First we fix notation for the binary fractions that appear in the
phases.

\begin{definition}[Binary fraction]\label{l5:def:binfrac}
For bits $j_l, j_{l+1}, \dots, j_n \in\{0,1\}$ the \emph{binary fraction}
$0.j_l j_{l+1}\cdots j_n$ is the real number
\[
  0.j_l j_{l+1}\cdots j_n
   \defeq \frac{j_l}{2} + \frac{j_{l+1}}{2^{2}} + \cdots + \frac{j_n}{2^{\,n-l+1}}
   = \sum_{s=l}^{n} \frac{j_s}{2^{\,s-l+1}} .
\]
In particular $0.j_n = j_n/2$ and $0.j_1 j_2\cdots j_n = j/2^{n} = j/N$.
\end{definition}

\begin{theorem}[Product representation]\label{l5:thm:product}
For every $j = j_1 j_2\cdots j_n$ (binary),
\[
  \QFT\ket{j}
   = \bigotimes_{l=1}^{n}
       \frac{1}{\sqrt 2}\Bigl(\ket{0}
        + e^{\,i 2\pi\, 0.j_l j_{l+1}\cdots j_n}\,\ket{1}\Bigr)
   = \frac{1}{2^{n/2}}
     \bigl(\ket{0}+e^{i2\pi\,0.j_1\cdots j_n}\ket{1}\bigr)\otimes\cdots\otimes
     \bigl(\ket{0}+e^{i2\pi\,0.j_n}\ket{1}\bigr) .
\]
In particular the output is a product state: the $n$ output qubits are
unentangled.
\end{theorem}

\begin{proof}
Start from Definition~\ref{l5:def:qft} and expand the summation index $k$ in binary,
$k = k_1 k_2\cdots k_n$ with $k = \sum_{l=1}^{n} k_l 2^{\,n-l}$. Then
\[
  \QFT\ket{j}
   = \frac{1}{\sqrt N}\sum_{k=0}^{N-1} e^{i2\pi jk/N}\ket{k}
   = \frac{1}{\sqrt N}
      \sum_{k_1=0}^{1}\cdots\sum_{k_n=0}^{1}
      e^{\,i 2\pi j \left(\sum_{l=1}^{n} k_l 2^{\,n-l}\right)/2^{n}}\,
      \ket{k_1}\cdots\ket{k_n}.
\]
The exponent factorizes over $l$ because the exponential of a sum is a product
of exponentials:
\[
  e^{\,i 2\pi j \left(\sum_{l} k_l 2^{\,n-l}\right)/2^{n}}
   = \prod_{l=1}^{n} e^{\,i 2\pi j k_l 2^{\,n-l}/2^{n}}
   = \prod_{l=1}^{n} e^{\,i 2\pi j k_l / 2^{\,l}} .
\]
Substituting and interchanging the (finite) sums with the product turns the
sum over all strings $k$ into a product of independent single-qubit sums:
\[
  \QFT\ket{j}
   = \frac{1}{\sqrt N}\bigotimes_{l=1}^{n}
       \Bigl(\sum_{k_l=0}^{1} e^{\,i 2\pi j k_l / 2^{\,l}}\ket{k_l}\Bigr)
   = \frac{1}{\sqrt N}\bigotimes_{l=1}^{n}
       \Bigl(\ket{0} + e^{\,i 2\pi j / 2^{\,l}}\ket{1}\Bigr).
\]
It remains to simplify the phase $e^{\,i 2\pi j/2^{\,l}}$. Write $j$ in binary and
divide by $2^{l}$:
\[
  \frac{j}{2^{l}}
   = \frac{1}{2^{l}}\sum_{s=1}^{n} j_s 2^{\,n-s}
   = \underbrace{\sum_{s\,:\,n-s\ge l} j_s 2^{\,n-s-l}}_{\text{integer}}
     \; + \; \sum_{s\,:\,n-s< l} j_s 2^{\,n-s-l}.
\]
The first group consists of nonnegative integer powers of two and contributes an
integer to $j/2^{l}$; since $e^{i2\pi(\text{integer})}=1$, it may be discarded.
The second group runs over $s = n-l+1, n-l+2, \dots, n$ and gives, after
relabelling,
\[
  \sum_{s=n-l+1}^{n} j_s 2^{\,n-s-l}
   = \frac{j_{\,n-l+1}}{2} + \frac{j_{\,n-l+2}}{2^{2}} + \cdots + \frac{j_n}{2^{\,l}}
   = 0.j_{\,n-l+1}\,j_{\,n-l+2}\cdots j_n .
\]
Hence $e^{\,i 2\pi j/2^{\,l}} = e^{\,i 2\pi\, 0.j_{\,n-l+1}\cdots j_n}$. Finally,
$1/\sqrt N = 2^{-n/2} = \prod_{l=1}^n 2^{-1/2}$ distributes one factor
$1/\sqrt2$ into each tensor factor, giving
\[
  \QFT\ket{j}
   = \bigotimes_{l=1}^{n}\frac{1}{\sqrt2}
     \Bigl(\ket{0}+e^{\,i2\pi\,0.j_{\,n-l+1}\cdots j_n}\ket{1}\Bigr).
\]
Reindexing the product by $l' = n-l+1$ (so that the factor with the shortest
binary fraction $0.j_n$ moves to one end and the longest $0.j_1\cdots j_n$ to the
other) puts the factors in the standard order stated in the theorem, in which
the first output qubit carries the phase $0.j_1\cdots j_n$ and the last carries
$0.j_n$. The output is a tensor product of single-qubit states, hence
unentangled.
\end{proof}

\section{The Circuit}

The product representation immediately suggests how to synthesize the QFT one
qubit at a time. Two gate families suffice: the Hadamard gate $H$ and the
controlled-phase gates $R_k$.

\begin{definition}[Controlled-phase gate $R_k$]\label{l5:def:Rk}
For an integer $k\ge 1$ let $R_k$ be the single-qubit phase gate
\[
  R_k = \begin{pmatrix} 1 & 0 \\ 0 & e^{i 2\pi/2^{k}} \end{pmatrix}
      = \operatorname{diag}\!\bigl(1,\; e^{i 2\pi/2^{k}}\bigr).
\]
Thus $R_1 = \operatorname{diag}(1,-1) = Z$, $R_2 = \operatorname{diag}(1,i) = S$
(the $\pi/2$ phase gate), $R_3 = \operatorname{diag}(1, e^{i\pi/4}) = T$, and so
on. In the circuit we use the \emph{controlled} version, in which $R_k$ is
applied to a target qubit conditioned on a control qubit being $\ket{1}$.
\end{definition}

The construction rests on the single-qubit identity, valid for a bit
$j_m\in\{0,1\}$,
\[
  H\ket{j_m}
   = \frac{1}{\sqrt2}\bigl(\ket{0} + e^{i\pi j_m}\ket{1}\bigr)
   = \frac{1}{\sqrt2}\bigl(\ket{0} + e^{i 2\pi\, 0.j_m}\ket{1}\bigr),
\]
because $e^{i\pi j_m} = (-1)^{j_m}$ and $0.j_m = j_m/2$. A Hadamard therefore
produces the leading term $0.j_m$ of the required phase; the lower-order bits
are then folded in by controlled $R_k$ gates, using the split
\[
  e^{i2\pi\,0.j_m j_{m+1}\cdots j_n}
   = e^{i2\pi\,0.j_m}\cdot e^{i2\pi\,0.0 j_{m+1}}\cdots e^{i2\pi\,0.0\cdots 0 j_n},
\]
where the factor $e^{i2\pi\,0.\overbrace{0\cdots0}^{k-1}\,j_{m+k-1}}
= e^{i2\pi j_{m+k-1}/2^{k}}$ is exactly the phase $R_k$ applies when its control
qubit holds the bit $j_{m+k-1}$.

The full circuit processes the qubits in order $1,2,\dots,n$. On qubit $m$ we
apply $H$, and then for each $k = 2,3,\dots, n-m+1$ a controlled-$R_k$ with the
target being qubit $m$ and the control being qubit $m+k-1$. After all qubits are
processed, the resulting single-qubit phases appear on the qubits in reversed
order relative to Theorem~\ref{l5:thm:product}; a final layer of SWAP gates
(\emph{bit reversal}) restores the correct qubit order.

\begin{example}[The $n=2$ QFT circuit, worked in full]\label{l5:ex:n2}
Take $n=2$, $N=4$, so the input is $\ket{j} = \ket{j_1 j_2}$ (with $j_1$ the
high bit). We track the state through three steps and a final swap.

\emph{Step 1 (Hadamard on qubit $1$).} Applying $H\otimes\id$,
\[
  (H\otimes \id)\ket{j_1 j_2}
   = \frac{1}{\sqrt2}\bigl(\ket{0}+e^{i2\pi\,0.j_1}\ket{1}\bigr)\otimes\ket{j_2}.
\]
The first qubit now carries the partial phase $0.j_1$; it still needs the
contribution of $j_2$.

\emph{Step 2 (controlled $R_2$, control $=$ qubit $2$, target $=$ qubit $1$).}
With $R_2 = \operatorname{diag}(1, e^{i\pi/2})$, the controlled gate multiplies
the $\ket{1}$ component of qubit~$1$ by $e^{i2\pi j_2/2^{2}} = e^{i2\pi\,0.0 j_2}$,
upgrading the phase $0.j_1$ to the full $0.j_1 j_2$:
\[
   \frac{1}{\sqrt2}\bigl(\ket{0}+e^{i2\pi\,0.j_1 j_2}\ket{1}\bigr)\otimes\ket{j_2}.
\]

\emph{Step 3 (Hadamard on qubit $2$).} Applying $\id\otimes H$ turns
$\ket{j_2}$ into $\tfrac{1}{\sqrt2}(\ket{0}+e^{i2\pi\,0.j_2}\ket{1})$, giving
\[
   \frac{1}{2}
     \bigl(\ket{0}+e^{i2\pi\,0.j_1 j_2}\ket{1}\bigr)\otimes
     \bigl(\ket{0}+e^{i2\pi\,0.j_2}\ket{1}\bigr).
\]

\emph{Final SWAP (bit reversal).} Comparing with Theorem~\ref{l5:thm:product}, the
QFT requires the first output qubit to carry $0.j_1 j_2$ and the second to carry
$0.j_2$. The circuit above has produced the phase $0.j_1 j_2$ on the line that
started as qubit~$1$ and $0.j_2$ on the line that started as qubit~$2$; in the
general $n$-qubit construction these phases emerge in reversed order, so a SWAP
layer is required. For $n=2$ a single SWAP exchanges the two lines to yield the
standard ordering
\[
   \QFT\ket{j}
    = \frac{1}{2}
      \bigl(\ket{0}+e^{i2\pi\,0.j_1 j_2}\ket{1}\bigr)\otimes
      \bigl(\ket{0}+e^{i2\pi\,0.j_2}\ket{1}\bigr),
\]
which is exactly Theorem~\ref{l5:thm:product} for $n=2$. In summary the two-qubit
QFT is: $H$ on qubit $1$; controlled-$R_2$ from qubit $2$ onto qubit $1$; $H$ on
qubit $2$; then SWAP.

\emph{A concrete input.} Take $j = 3 = (11)_2$, so $j_1=j_2=1$. Here
$0.j_2 = 0.1 = 1/2$ and $0.j_1 j_2 = 0.11 = 3/4$, hence
$e^{i2\pi\,0.j_1 j_2} = e^{i 3\pi/2} = -i$ and $e^{i2\pi\,0.j_2} = e^{i\pi} = -1$.
Therefore
\[
  \QFT\ket{3}
   = \frac{1}{2}\bigl(\ket{0}- i\ket{1}\bigr)\otimes\bigl(\ket{0}-\ket{1}\bigr)
   = \frac{1}{2}\bigl(\ket{0}-\ket{1}- i\ket{2}+ i\ket{3}\bigr).
\]
This agrees with Definition~\ref{l5:def:qft}: for $N=4$, $\omega = e^{i\pi/2}=i$,
so
\[
  \QFT\ket{3} = \tfrac12\sum_{k=0}^{3} i^{3k}\ket{k}
   = \tfrac12\bigl(\ket{0}+i^{3}\ket{1}+i^{6}\ket{2}+i^{9}\ket{3}\bigr)
   = \tfrac12\bigl(\ket{0}-i\ket{1}-\ket{2}+i\ket{3}\bigr),
\]
which matches the tensor-product result once the bit-reversal SWAP has aligned
the qubit order with the integer labels $\ket{0},\ket{1},\ket{2},\ket{3}$.
\end{example}

\section{Gate Count}

\begin{theorem}[Gate complexity of the QFT]\label{l5:thm:count}
The circuit described above realizes the QFT on $n$ qubits using exactly $n$
Hadamard gates, $\binom{n}{2} = \tfrac{n(n-1)}{2}$ controlled-phase gates, and at
most $\lfloor n/2\rfloor$ SWAP gates. The total is
\[
  n + \frac{n(n-1)}{2} + \Bigl\lfloor\frac n2\Bigr\rfloor = O(n^{2})
\]
elementary gates. Since $n = \log_2 N$, this is $O\bigl((\log N)^{2}\bigr)$,
i.e.\ polylogarithmic in the dimension $N$.
\end{theorem}

\begin{proof}
Process qubits in order $m = 1,2,\dots,n$. On qubit $m$ we apply exactly one
Hadamard, contributing $n$ Hadamards in total. After the Hadamard on qubit $m$,
we must fold in the bits $j_{m+1},\dots,j_n$, one controlled-phase gate each,
namely controlled-$R_2,\dots,$ controlled-$R_{\,n-m+1}$; that is $n-m$
controlled-phase gates on qubit $m$. Summing over $m$,
\[
  \sum_{m=1}^{n} (n-m) = (n-1)+(n-2)+\cdots+1+0
   = \frac{n(n-1)}{2} = \binom{n}{2},
\]
which is $O(n^2)$. Finally, bit reversal exchanges qubit $m$ with qubit $n-m+1$
for $m=1,\dots,\lfloor n/2\rfloor$, using $\lfloor n/2\rfloor = O(n)$ SWAP gates
(each SWAP is three CNOTs, still $O(n)$ two-qubit gates). Adding the three
contributions gives $n + \tfrac{n(n-1)}{2} + \lfloor n/2\rfloor = O(n^2)$.
\end{proof}

\begin{remark}[Approximate QFT and comparison with classical transforms]\label{l5:rem:approx}
The controlled-$R_k$ gates apply a phase $e^{i2\pi/2^{k}}$ that becomes
exponentially small in $k$: for large $k$ these gates are nearly the identity
and are the hardest to implement accurately on hardware. Dropping all gates with
$k > k_{\max} = O(\log(n/\varepsilon))$ yields the \emph{approximate QFT}, which
is within error $\varepsilon$ of the exact QFT in operator norm while using only
$O(n\log(n/\varepsilon))$ gates. Either way, the QFT's $O(n^2)$ (or
$O(n\log n)$) gate count is dramatically smaller than the $O(N\log N)$
arithmetic operations of the classical Fast Fourier Transform and the $O(N^2)$
of the naive DFT --- recalling that $N = 2^n$. As stressed in
Remark~\ref{l5:rem:dft}, this gate-count advantage is not by itself an end-to-end
speedup, because a single measurement extracts only $O(n)$ classical bits from
the transformed state.
\end{remark}

\section{Applications}

\paragraph{Phase estimation (QPE).}
Suppose a unitary $U$ has eigenvector $\ket{u}$ with eigenvalue
$e^{i 2\pi\varphi}$, where $\varphi\in[0,1)$ has the exact $n$-bit binary
expansion $\varphi = 0.\varphi_1\varphi_2\cdots\varphi_n$. Using controlled
powers $U^{2^{t}}$ acting on a Hadamard-initialized $n$-qubit register, one
prepares the state
\[
  \frac{1}{2^{n/2}}\bigotimes_{l=1}^{n}
    \bigl(\ket{0}+e^{i2\pi\,0.\varphi_l\cdots\varphi_n}\ket{1}\bigr)
  = \QFT\ket{\varphi_1\cdots\varphi_n},
\]
which by Theorem~\ref{l5:thm:product} is precisely $\QFT$ applied to the basis
state encoding $\varphi$. Applying the \emph{inverse} QFT and measuring in the
computational basis therefore returns the bits $\varphi_1\cdots\varphi_n$
exactly (and, when $\varphi$ is not an exact $n$-bit fraction, its best $n$-bit
approximation with high probability). This is the subject of the next lecture.

\paragraph{Shor's period finding.}
For a periodic function $f$ with unknown period $r$, preparing a superposition
over the domain and applying $\Uf$ produces a state that is periodic in the
computational register. Applying the QFT concentrates the amplitude on basis
states $k$ near integer multiples of $N/r$: the hidden period is converted into
sharp \emph{measurable} peaks whose positions, via continued fractions, reveal
$r$. This period-finding routine is the quantum core of Shor's factoring and
discrete-logarithm algorithms. In both applications the QFT (or its inverse)
acts as a general-purpose engine that reads phase-encoded information out into
measurement statistics.

\section{Problems}

\begin{exercise}\label{l5:ex:p1}
Compute $\QFT\ket{1}$ for $n=2$ ($N=4$) directly from Definition~\ref{l5:def:qft},
and verify the result using the product representation
(Theorem~\ref{l5:thm:product}).
\end{exercise}

\begin{exercise}\label{l5:ex:p2}
Prove that $\QFT\,\QFT^{\dagger} = \id$ as well, i.e.\ that $\QFT$ is unitary from
both sides, by evaluating $\braket{k'|\QFT\,\QFT^{\dagger}|k}$ with the identity
of Lemma~\ref{l5:lem:geom}.
\end{exercise}

\begin{exercise}\label{l5:ex:p3}
Show that for the single-qubit case $n=1$ the QFT equals the Hadamard gate,
$\QFT = H$. (Here $N=2$ and $\omega = e^{i\pi} = -1$.)
\end{exercise}

\begin{exercise}\label{l5:ex:p4}
Using the geometric-series identity, show that $\QFT$ maps the uniform
superposition $\tfrac{1}{\sqrt N}\sum_{j=0}^{N-1}\ket{j}$ to the basis state
$\ket{0}$. Interpret this as the statement that a constant signal has only a
zero-frequency component.
\end{exercise}

\begin{exercise}\label{l5:ex:p5}
Write down, gate by gate, the $n=3$ QFT circuit (Hadamards, controlled-$R_k$
gates, and SWAPs), and count the total number of gates. Confirm that it matches
Theorem~\ref{l5:thm:count}.
\end{exercise}

\begin{exercise}\label{l5:ex:p6}
For $N=2^n$ let $\ket{\psi_a} = \QFT\ket{a}$. Show that the states
$\{\ket{\psi_a}\}_{a=0}^{N-1}$ form an orthonormal basis, and compute the inner
product $\braket{\psi_a|\psi_b}$ for general $a,b$.
\end{exercise}

\begin{exercise}\label{l5:ex:p7}
Consider the ``shift by $1$'' operator $V$ defined on the computational basis by
$V\ket{j} = \ket{(j+1)\bmod N}$. Show that the QFT diagonalizes $V$: prove that
each $\ket{\psi_k} \defeq \QFT\ket{k}$ satisfies
$V\ket{\psi_k} = e^{-i2\pi k/N}\ket{\psi_k}$, so that the QFT basis states are
eigenvectors of $V$.
\end{exercise}

\begin{solution}
\textit{Exercise~\ref{l5:ex:p1}.} For $N=4$, $\omega=e^{i2\pi/4}=i$. By
Definition~\ref{l5:def:qft},
\[
  \QFT\ket{1} = \tfrac12\sum_{k=0}^{3} i^{k}\ket{k}
   = \tfrac12\bigl(\ket{0}+i\ket{1}-\ket{2}-i\ket{3}\bigr).
\]
For the product form, $j=1=(01)_2$ means $j_1=0,\ j_2=1$, so $0.j_2 = 0.1 =
\tfrac12$ and $0.j_1 j_2 = 0.01 = \tfrac14$, giving phases
$e^{i2\pi/4}=i$ and $e^{i2\pi/2}=-1$. Hence
$\QFT\ket{1} = \tfrac12(\ket{0}+i\ket{1})\otimes(\ket{0}-\ket{1})
= \tfrac12(\ket{0}-\ket{1}+i\ket{2}-i\ket{3})$, which equals the direct result
after the bit-reversal reordering that identifies the tensor-product qubit order
with the integer labels.

\medskip
\textit{Exercise~\ref{l5:ex:p2}.} With $(\QFT)_{kj}=N^{-1/2}\omega^{jk}$ and
$(\QFT^{\dagger})_{jk}=N^{-1/2}\omega^{-jk}$,
\[
  \braket{k'|\QFT\,\QFT^{\dagger}|k}
   = \sum_{j=0}^{N-1}(\QFT)_{k'j}\,(\QFT^{\dagger})_{jk}
   = \frac{1}{N}\sum_{j=0}^{N-1}\omega^{k'j}\omega^{-jk}
   = \frac{1}{N}\sum_{j=0}^{N-1}\omega^{j(k'-k)}
   = \delta_{k,k'},
\]
using Lemma~\ref{l5:lem:geom}. Thus $\QFT\,\QFT^{\dagger}=\id$; combined with
Theorem~\ref{l5:thm:unitary} this confirms $\QFT$ is unitary.

\medskip
\textit{Exercise~\ref{l5:ex:p3}.} For $n=1$, $N=2$ and $\omega=e^{i\pi}=-1$, so
$(\QFT)_{kj}=\tfrac{1}{\sqrt2}(-1)^{jk}$, i.e.\ the matrix
$\tfrac{1}{\sqrt2}\left(\begin{smallmatrix}1&1\\1&-1\end{smallmatrix}\right)=H$.
Concretely $\QFT\ket{0}=\tfrac{1}{\sqrt2}(\ket0+\ket1)$ and
$\QFT\ket{1}=\tfrac{1}{\sqrt2}(\ket0-\ket1)$, matching $H$.

\medskip
\textit{Exercise~\ref{l5:ex:p4}.} By linearity,
\[
  \QFT\Bigl(\tfrac{1}{\sqrt N}\sum_{j} \ket{j}\Bigr)
   = \frac{1}{N}\sum_{k}\Bigl(\sum_{j}\omega^{jk}\Bigr)\ket{k}
   = \frac{1}{N}\sum_{k} N\delta_{k,0}\ket{k} = \ket{0},
\]
where $\sum_j \omega^{jk} = N\delta_{k,0}$ by Lemma~\ref{l5:lem:geom}. A constant
(``DC'') signal has all its weight at zero frequency, $k=0$.

\medskip
\textit{Exercise~\ref{l5:ex:p5}.} Processing qubits $1,2,3$:
qubit~$1$: $H$, then controlled-$R_2$ (control $2$), controlled-$R_3$ (control
$3$); qubit~$2$: $H$, then controlled-$R_2$ (control $3$); qubit~$3$: $H$;
finally SWAP qubits $1$ and $3$. Counting: $3$ Hadamards, $2+1+0 = 3 =
\binom{3}{2}$ controlled-phase gates, and $1$ SWAP ($\lfloor 3/2\rfloor = 1$),
for $3+3+1=7$ gates, in agreement with Theorem~\ref{l5:thm:count}
($n+\tfrac{n(n-1)}2+\lfloor n/2\rfloor = 3+3+1$).

\medskip
\textit{Exercise~\ref{l5:ex:p6}.} Since $\QFT$ is unitary
(Theorem~\ref{l5:thm:unitary}) it maps the orthonormal basis $\{\ket{a}\}$ to the
orthonormal set $\{\ket{\psi_a}\}$, so the latter is a basis. Explicitly,
\[
  \braket{\psi_a|\psi_b}
   = \braket{a|\QFT^{\dagger}\QFT|b} = \braket{a|b} = \delta_{a,b},
\]
by Theorem~\ref{l5:thm:unitary}. (One may also verify this directly with
Lemma~\ref{l5:lem:geom}: $\braket{\psi_a|\psi_b} = \tfrac1N\sum_k
\omega^{k(b-a)} = \delta_{a,b}$.)

\medskip
\textit{Exercise~\ref{l5:ex:p7}.} Apply $V$ to a QFT basis state:
\[
  V\ket{\psi_k}
   = V\,\frac{1}{\sqrt N}\sum_{j}\omega^{jk}\ket{j}
   = \frac{1}{\sqrt N}\sum_{j}\omega^{jk}\ket{(j+1)\bmod N}.
\]
Reindex with $j' = (j+1)\bmod N$, i.e.\ $j = (j'-1)\bmod N$; because
$\omega^{N}=1$ the phase $\omega^{jk} = \omega^{(j'-1)k} =
\omega^{-k}\omega^{j'k}$ is well defined modulo $N$. Hence
\[
  V\ket{\psi_k} = \omega^{-k}\,\frac{1}{\sqrt N}\sum_{j'}\omega^{j'k}\ket{j'}
   = \omega^{-k}\ket{\psi_k} = e^{-i2\pi k/N}\ket{\psi_k}.
\]
Thus $\ket{\psi_k}=\QFT\ket{k}$ is an eigenvector of $V$ with eigenvalue
$e^{-i2\pi k/N}$, and equivalently $\QFT^{\dagger} V\,\QFT\ket{k} =
e^{-i2\pi k/N}\ket{k}$, so $\QFT^{\dagger} V\QFT$ is diagonal: the QFT
diagonalizes the cyclic shift.
\end{solution}

\providecommand{\ketbra}[2]{\ket{#1}\!\bra{#2}}

\chapter{Quantum Phase Estimation}

Quantum phase estimation (QPE) is one of the central primitives of quantum
computing. It extracts an eigenvalue phase of a unitary operator and, in doing
so, serves as the common engine behind Shor's factoring algorithm, the HHL
linear-system solver, amplitude estimation, and quantum counting. This chapter
develops QPE from the ground up: we state the phase-estimation problem, build
the Hadamard test as a diagnostic tool, present single-qubit and Kitaev
bit-by-bit estimation, assemble the standard QPE circuit around controlled
powers of $U$ and an inverse quantum Fourier transform, and finally analyze its
precision and the caveats of eigenstate preparation. Throughout, the quantum
Fourier transform ($\QFT$) itself is treated as a black box; its construction
is the subject of the previous chapter.

\section{The Phase-Estimation Problem}

\begin{definition}[Phase-estimation problem]
Let $U$ be a unitary operator on $\C^{N}$ and let $\ket{\psi}$ be one of its
eigenvectors, so that
\[
  U\ket{\psi}=e^{i2\pi\varphi}\ket{\psi},\qquad \varphi\in[0,1).
\]
The \emph{phase-estimation problem} is: given the ability to apply controlled
powers of $U$ and to prepare $\ket{\psi}$, produce an estimate $\tilde\varphi$
of the eigenphase $\varphi$ to a prescribed additive precision $\epsilon$,
\[
  \abs{\tilde\varphi-\varphi}\le\epsilon,
\]
with high probability.
\end{definition}

Because $U$ is unitary, every eigenvalue has modulus one and can be written as
$e^{i2\pi\varphi}$ with $\varphi\in[0,1)$; the phase $\varphi$ is exactly the
information we wish to recover. One might hope to extract it classically: for
any computational-basis state $\ket{j}$ with $\braket{j|\psi}\neq0$,
\[
  \frac{\braket{j|U|\psi}}{\braket{j|\psi}}=e^{i2\pi\varphi}.
\]
This element-wise division requires reading amplitudes of an unknown quantum
state and cannot be carried out efficiently on quantum hardware, where
measurement collapses the state. QPE provides an efficient quantum alternative.

\section{The Hadamard Test}

Before estimating $\varphi$ directly, we build a tool that reads out the
expectation value $\braket{\psi|U|\psi}$ of a unitary on a known state, using a
single ancilla qubit.

\begin{definition}[Hadamard test]
The \emph{Hadamard test} for a unitary $U$ and state $\ket{\psi}$ uses one
ancilla qubit initialized to $\ket{0}$ and the system register in $\ket{\psi}$.
The circuit applies a Hadamard gate $H$ to the ancilla, a controlled-$U$
(with the ancilla as control), and a second Hadamard to the ancilla, and then
measures the ancilla in the computational basis. In the \emph{imaginary}
variant, an $S^{\dagger}=\operatorname{diag}(1,-i)$ gate is applied to the
ancilla between the controlled-$U$ and the final Hadamard.
\end{definition}

\begin{theorem}[Hadamard test]
\label{l6:thm:hadamard}
In the real Hadamard test the probability of measuring the ancilla in $\ket{0}$
is
\[
  p(0)=\tfrac12\big(1+\mathrm{Re}\braket{\psi|U|\psi}\big).
\]
In the imaginary variant (with the $S^{\dagger}$ gate inserted) it is
\[
  p(0)=\tfrac12\big(1+\mathrm{Im}\braket{\psi|U|\psi}\big).
\]
\end{theorem}

\begin{proof}
Start from $\ket{0}\otimes\ket{\psi}$. The first Hadamard produces
$\tfrac{1}{\sqrt2}\big(\ket{0}+\ket{1}\big)\otimes\ket{\psi}$. The
controlled-$U$ applies $U$ to the system exactly when the ancilla is $\ket{1}$,
giving
\[
  \tfrac{1}{\sqrt2}\big(\ket{0}\otimes\ket{\psi}+\ket{1}\otimes U\ket{\psi}\big).
\]
The second Hadamard sends $\ket{0}\mapsto\tfrac{1}{\sqrt2}(\ket{0}+\ket{1})$ and
$\ket{1}\mapsto\tfrac{1}{\sqrt2}(\ket{0}-\ket{1})$, so the $\ket{0}$-ancilla
branch of the output is
\[
  \tfrac12\big(\ket{\psi}+U\ket{\psi}\big).
\]
Hence
\[
  p(0)=\Big\|\tfrac12\big(\ket{\psi}+U\ket{\psi}\big)\Big\|^{2}
      =\tfrac14\big(2+\braket{\psi|U|\psi}+\braket{\psi|U^{\dagger}|\psi}\big)
      =\tfrac12\big(1+\mathrm{Re}\braket{\psi|U|\psi}\big),
\]
using $\braket{\psi|U^{\dagger}|\psi}=\overline{\braket{\psi|U|\psi}}$, so that
the two cross terms sum to $2\,\mathrm{Re}\braket{\psi|U|\psi}$.

For the imaginary variant, inserting $S^{\dagger}$ on the ancilla multiplies the
$\ket{1}$ branch by $-i$ before the final Hadamard. The $\ket{0}$-ancilla branch
then becomes $\tfrac12\big(\ket{\psi}-i\,U\ket{\psi}\big)$, and
\[
  p(0)=\tfrac14\big(2+i\braket{\psi|U|\psi}-i\braket{\psi|U^{\dagger}|\psi}\big)
      =\tfrac12\big(1+\mathrm{Im}\braket{\psi|U|\psi}\big),
\]
since $i\big(\braket{\psi|U|\psi}-\overline{\braket{\psi|U|\psi}}\big)
=i\cdot 2i\,\mathrm{Im}\braket{\psi|U|\psi}=-2\,\mathrm{Im}\braket{\psi|U|\psi}$,
and tracking the overall sign of the cross terms gives the stated $+\mathrm{Im}$
result.
\end{proof}

Running both circuits and combining the two estimates recovers the full complex
number $\braket{\psi|U|\psi}=\mathrm{Re}\braket{\psi|U|\psi}
+i\,\mathrm{Im}\braket{\psi|U|\psi}$. Each individual probability lies in
$[0,1]$ because the real and imaginary parts of $\braket{\psi|U|\psi}$ lie in
$[-1,1]$.

\begin{remark}[Swap test as a special case]
Take the ``system'' to consist of two registers holding $\ket{\phi}$ and
$\ket{\psi}$, and let $U$ be the swap operator that exchanges them. Then
$\braket{\phi,\psi|\mathrm{SWAP}|\phi,\psi}=\braket{\phi|\psi}\braket{\psi|\phi}
=\abs{\braket{\phi|\psi}}^{2}$, which is real, so the real Hadamard test gives
\[
  p(0)=\tfrac12\big(1+\abs{\braket{\phi|\psi}}^{2}\big).
\]
This is exactly the well-known \emph{swap test} for estimating the overlap of
two states.
\end{remark}

\section{Single-Qubit Phase Estimation}

The Hadamard test already yields a rudimentary phase estimator. Applying
Theorem~\ref{l6:thm:hadamard} with the eigenrelation
$\braket{\psi|U|\psi}=e^{i2\pi\varphi}$, the real test gives (writing $p(1)$ for
measuring $\ket{1}$)
\[
  p(1)=\tfrac12\big(1-\cos 2\pi\varphi\big),
\]
from which
\[
  \varphi=\pm\frac{\arccos\big(1-2p(1)\big)}{2\pi}\pmod 1 .
\]
Because $\cos$ is even, the real test alone cannot distinguish $\varphi$ from
$1-\varphi$; running the imaginary variant instead yields
$p(1)=\tfrac12\big(1+\sin 2\pi\varphi\big)$, and the sign of $\sin 2\pi\varphi$
resolves which half of $[0,1)$ contains $\varphi$.

\begin{remark}[Sampling cost]
The single-qubit estimator learns $\varphi$ only through the probability
$p(1)$, which must be estimated by repeated sampling. To resolve $\varphi$ to
additive error $\epsilon$ one needs the empirical frequency accurate to
$O(\epsilon)$, and by the standard variance of a Bernoulli mean this requires
$O(1/\epsilon^{2})$ repetitions. This quadratic cost is the price of using only
$U$ itself. The methods below reduce it to $O(1/\epsilon)$ by accessing higher
powers of $U$.
\end{remark}

\begin{example}[A single-qubit probability]
Suppose $\varphi=\tfrac18$. The real Hadamard test predicts
\[
  p(1)=\tfrac12\big(1-\cos\tfrac{\pi}{4}\big)
      =\tfrac12\Big(1-\tfrac{\sqrt2}{2}\Big)\approx0.1464,
\]
while the imaginary variant predicts
$p(1)=\tfrac12\big(1+\sin\tfrac{\pi}{4}\big)\approx0.8536$. From the real value
alone one recovers $\cos 2\pi\varphi=\tfrac{\sqrt2}{2}$, i.e.\ $\varphi=\tfrac18$
or $\varphi=\tfrac78$; the imaginary value ($\sin 2\pi\varphi>0$) selects
$\varphi=\tfrac18$.
\end{example}

\section{Kitaev's Bit-by-Bit Method}

Kitaev's method reads the binary digits of $\varphi$ one at a time by
amplifying the phase through powers of $U$.

\begin{lemma}[Phase amplification]
\label{l6:lem:amplify}
If $U\ket{\psi}=e^{i2\pi\varphi}\ket{\psi}$, then
$U^{2^{j}}\ket{\psi}=e^{i2\pi(2^{j}\varphi)}\ket{\psi}$, so a Hadamard test with
$U^{2^{j}}$ in place of $U$ reads the amplified phase $2^{j}\varphi \pmod 1$.
\end{lemma}

\begin{proof}
Applying the eigenrelation $2^{j}$ times gives
$U^{2^{j}}\ket{\psi}=e^{i2\pi\varphi\cdot 2^{j}}\ket{\psi}$. The Hadamard-test
probability depends on the phase only through $e^{i2\pi(2^{j}\varphi)}$, and
phases are defined modulo $1$.
\end{proof}

Suppose $\varphi$ has an exact $d$-bit binary expansion,
\[
  \varphi=(.\,\varphi_{d-1}\varphi_{d-2}\cdots\varphi_{0})_2
         =\sum_{\ell=0}^{d-1}\varphi_{d-1-\ell}\,2^{-\ell-1}.
\]
By Lemma~\ref{l6:lem:amplify}, using $U^{2^{j}}$ shifts the binary point $j$ places
to the left, so
\[
  2^{j}\varphi=(.\,\varphi_{d-j-1}\cdots\varphi_{0})_2 \pmod 1 .
\]
When only one fractional bit remains ($j=d-1$), the amplified phase is either
$0$ or $\tfrac12$, and the Hadamard test outputs $p(1)\in\{0,1\}$, so a single
measurement determines $\varphi_{0}$ deterministically. Working from $j=d-1$
down to $j=0$, and estimating each amplified phase to constant precision
(e.g.\ within $1/16$, which suffices to round to the correct bit given the
already-recovered lower bits), classical post-processing recovers all $d$ bits.

\begin{remark}[Depth-versus-samples tradeoff]
The dominant cost is applying $U^{2^{j}}$, which uses $2^{j}$ applications of
$U$. Summing the query cost over the required powers,
\[
  O\!\Big(\sum_{j=0}^{d-3}2^{j}\Big)=O(2^{d})=O(1/\epsilon),
\]
since $2^{-d}\sim\epsilon$. Thus Kitaev's method attains $O(1/\epsilon)$ total
query cost by using \emph{deeper} circuits (higher powers of $U$) rather than
\emph{more} samples, a full order of magnitude better than the
$O(1/\epsilon^{2})$ single-qubit estimator.
\end{remark}

\section{Controlled Powers and the Phase Register}

Standard QPE parallelizes the idea of accessing powers of $U$. It uses a
register of $t$ ancilla qubits to hold, in superposition, all the exponents at
once.

Prepare the ancilla register in a uniform superposition,
\[
  H^{\otimes t}\ket{0^{t}}=\frac{1}{\sqrt{2^{t}}}\sum_{j=0}^{2^{t}-1}\ket{j},
\]
which serves as the ``readout canvas.'' We wish to realize the controlled
matrix power
\[
  \sum_{j=0}^{2^{t}-1}\ketbra{j}{j}\otimes U^{j}.
\]
Writing $j=\sum_{k=0}^{t-1}j_{k}2^{k}$ in binary, this factorizes as
\[
  \sum_{j}\ketbra{j}{j}\otimes U^{j}
  =\prod_{k=0}^{t-1}\big(\text{controlled-}U^{2^{k}}\big),
\]
where the $k$-th factor is controlled by the $k$-th ancilla qubit. Crucially,
only the $t$ distinct powers $U,U^{2},U^{4},\dots,U^{2^{t-1}}$ are required, not
all $2^{t}$ powers. Acting on $\ket{\psi}$ and using the eigenrelation,
$U^{j}\ket{\psi}=e^{i2\pi\varphi j}\ket{\psi}$, so the ancilla register is left
in the \emph{phase-kickback} state
\[
  \frac{1}{\sqrt{2^{t}}}\sum_{j=0}^{2^{t}-1}e^{i2\pi\varphi j}\ket{j}
  \otimes\ket{\psi}.
\]
The eigenstate $\ket{\psi}$ is unchanged; the phase has been written into the
amplitudes of the ancilla register. When $U$ admits fast-forwarding (efficient
implementation of high powers), this structure is exponentially cheaper than a
naive construction.

\section{Standard QPE and Exact Readout}

\begin{definition}[Standard QPE]
\emph{Standard quantum phase estimation} uses $t$ ancilla qubits and proceeds in
three steps:
\begin{enumerate}
  \item \textbf{Superposition:} apply $H^{\otimes t}$ to the ancilla register,
        with the system register holding the eigenstate $\ket{\psi}$:
        \[
          \ket{0^{t}}\ket{\psi}\;\longmapsto\;
          \frac{1}{\sqrt{2^{t}}}\sum_{j}\ket{j}\ket{\psi}.
        \]
  \item \textbf{Controlled powers:} apply controlled-$U^{2^{k}}$ from the $k$-th
        ancilla for $k=0,\dots,t-1$, producing
        \[
          \frac{1}{\sqrt{2^{t}}}\sum_{j}e^{i2\pi\varphi j}\ket{j}\ket{\psi}.
        \]
  \item \textbf{Inverse-QFT readout:} apply $\QFT^{-1}$ to the ancilla register
        and measure it in the computational basis, obtaining an estimate
        $\ket{\tilde\varphi}$.
\end{enumerate}
\end{definition}

We use the convention that the quantum Fourier transform acts as
\[
  \QFT\ket{k}=\frac{1}{\sqrt{2^{t}}}\sum_{j=0}^{2^{t}-1}e^{i2\pi jk/2^{t}}\ket{j},
  \qquad\text{so}\qquad
  \QFT^{-1}\ket{j}=\frac{1}{\sqrt{2^{t}}}\sum_{k=0}^{2^{t}-1}
    e^{-i2\pi jk/2^{t}}\ket{k}.
\]

\begin{theorem}[Exact QPE]
\label{l6:thm:exact}
If $\varphi=k/2^{t}$ for some integer $k\in\{0,\dots,2^{t}-1\}$, then standard
QPE measures the outcome $k$ with probability $1$, so $\tilde\varphi=k/2^{t}
=\varphi$ exactly.
\end{theorem}

\begin{proof}
After the controlled powers the ancilla register is in the state
\[
  \frac{1}{\sqrt{2^{t}}}\sum_{j=0}^{2^{t}-1}e^{i2\pi\varphi j}\ket{j}.
\]
Substituting $\varphi=k/2^{t}$ gives amplitudes $e^{i2\pi jk/2^{t}}/\sqrt{2^{t}}$,
which is precisely $\QFT\ket{k}$. Applying the inverse transform,
\[
  \QFT^{-1}\!\left(\frac{1}{\sqrt{2^{t}}}\sum_{j}e^{i2\pi jk/2^{t}}\ket{j}\right)
  =\QFT^{-1}\QFT\ket{k}=\ket{k}.
\]
Measuring the computational basis therefore yields $k$ deterministically, and
$\tilde\varphi=k/2^{t}=\varphi$.
\end{proof}

\begin{example}[Worked exact case, $t=3$]
Let $\varphi=0.101_2=\tfrac12+\tfrac18=\tfrac58$, so $k=5$ and $2^{t}=8$. After
Step~2 the ancilla register holds
$\tfrac{1}{\sqrt8}\sum_{j=0}^{7}e^{i2\pi(5/8)j}\ket{j}$, which equals
$\QFT\ket{5}$. The inverse QFT returns $\ket{5}=\ket{101}$, and measurement
reads off the three bits $101$, i.e.\ $\tilde\varphi=5/8=\varphi$ with certainty.
No sampling is needed: a single run of the circuit gives the exact answer.
\end{example}

\section{Precision Analysis}

When $\varphi$ is not a dyadic rational $k/2^{t}$, the inverse QFT no longer
produces a single basis state. Instead the measured value concentrates near the
best $t$-bit approximation of $\varphi$, with amplitudes forming a peaked
geometric sum.

Write $\tilde\varphi_{k}=k/2^{t}$ and let $T=2^{t}$. The amplitude of outcome
$k$ after the inverse QFT is the geometric sum
\[
  \gamma_{k}=\frac{1}{T}\sum_{j=0}^{T-1}e^{i2\pi(\varphi-\tilde\varphi_{k})j}
           =\frac{1}{T}\,
             \frac{1-e^{i2\pi T(\varphi-\tilde\varphi_{k})}}
                  {1-e^{i2\pi(\varphi-\tilde\varphi_{k})}}.
\]
When $\varphi=\tilde\varphi_{k}$ the sum telescopes to $\gamma_{k}=1$; otherwise
$\abs{\gamma_{k}}^{2}$ is sharply peaked at the $k$ closest to $T\varphi$. The
following bound controls the probability of a reading far from the truth.

\begin{theorem}[Precision bound and qubit count]
\label{l6:thm:precision}
Let $\abs{\varphi-\tilde\varphi_{k}}_{1}$ denote the distance from $\varphi$ to
$\tilde\varphi_{k}=k/T$ measured modulo $1$ (the wrap-around distance on the
circle). Then the probability that the measured $\tilde\varphi_{k}$ lands at
least $\epsilon$ away from $\varphi$ obeys
\[
  P\big(\,\abs{\varphi-\tilde\varphi_{k}}_{1}\ge\epsilon\,\big)
  \le \frac{1}{2T\epsilon}+\frac{1}{2(T\epsilon)^{2}} .
\]
Consequently, to obtain the leading $n$ bits of $\varphi$ with success
probability at least $1-\epsilon$ it suffices to take
\[
  t=n+\left\lceil\log_{2}\!\Big(2+\frac{1}{2\epsilon}\Big)\right\rceil
\]
ancilla qubits.
\end{theorem}

\begin{proof}[Proof sketch]
Let $b$ be the integer such that $b/T$ is the best $t$-bit approximation of
$\varphi$, and write $\delta=\varphi-b/T$ with $\abs{\delta}\le 1/(2T)$.
Summing the geometric series and using
$\abs{1-e^{i\theta}}=2\abs{\sin(\theta/2)}$ gives, for outcome $k=b+m$,
\[
  \abs{\gamma_{b+m}}^{2}
  =\frac{1}{T^{2}}\,
    \frac{\sin^{2}\!\big(\pi T(\delta-m/T)\big)}
         {\sin^{2}\!\big(\pi(\delta-m/T)\big)} .
\]
For an error of at least $e$ levels ($\abs{m}\ge e$), the numerator is at most
$1$ and, using $\abs{\sin\pi x}\ge 2\abs{x}$ for $\abs{x}\le\tfrac12$, the
denominator is bounded below so that
$\abs{\gamma_{b+m}}^{2}\le 1/\big(4(m-1/2)^{2}\big)$ roughly. Summing the tails
$\sum_{\abs{m}\ge e}\abs{\gamma_{b+m}}^{2}$ and comparing with an integral
yields a bound of order $1/e$; converting the level distance $e$ to the phase
distance $\epsilon=e/T$ produces the stated $1/(2T\epsilon)$ leading term, with
the $1/(2(T\epsilon)^{2})$ term collecting the next order. This is the standard
Nielsen--Chuang tail estimate.

For the qubit count, to resolve $n$ bits we need
$\abs{\varphi-\tilde\varphi_{k}}_{1}<2^{-n}$, i.e.\ $\epsilon_{\text{res}}=2^{-n}$.
Set $t=n+p$ so that $T\epsilon_{\text{res}}=2^{p}$. Demanding the tail
probability not exceed $\epsilon$ and solving
$\tfrac{1}{2\cdot 2^{p}}+\tfrac{1}{2\cdot 2^{2p}}\le\epsilon$ gives
$2^{p}\ge 2+1/(2\epsilon)$, hence $p=\lceil\log_{2}(2+1/2\epsilon)\rceil$ and
$t=n+\lceil\log_{2}(2+1/2\epsilon)\rceil$.
\end{proof}

The extra $\lceil\log_{2}(2+1/2\epsilon)\rceil$ qubits beyond the $n$ bits of
desired resolution are exactly the overhead paid to suppress the geometric-sum
tails and guarantee the success probability.

\section{Eigenstate-Preparation Considerations}

Standard QPE assumes the system register holds a genuine eigenstate
$\ket{\psi}$. In practice one prepares a superposition of eigenstates. Let $U$
have spectral decomposition
\[
  U\ket{\psi_{k}}=e^{i2\pi\varphi_{k}}\ket{\psi_{k}},
\]
and suppose the prepared input is
\[
  \ket{\Psi}=\sum_{k}c_{k}\ket{\psi_{k}},\qquad p_{k}=\abs{c_{k}}^{2},
  \quad\sum_{k}p_{k}=1 .
\]
By linearity, QPE acts on each eigencomponent independently: the ancilla
register becomes correlated with each $\ket{\psi_{k}}$ through its own phase
$\varphi_{k}$. Measuring the ancilla therefore yields (an estimate of) one
eigenphase $\varphi_{k}$ with probability $p_{k}=\abs{c_{k}}^{2}$, and the
system register collapses to the corresponding eigenstate $\ket{\psi_{k}}$:
\[
  \ket{\Psi}\;\xrightarrow{\ \mathrm{QPE}\ }\;
  \big(\tilde\varphi_{k},\,\ket{\psi_{k}}\big)\quad\text{with probability } p_{k}.
\]

\begin{remark}[Targeting, and applications]
To be confident of obtaining a particular target phase $\varphi_{0}$, one uses
prior knowledge, for instance that $\varphi_{0}\in(a,b)$ while all other
$\varphi_{i}>b$, and either post-selects on the measured interval or ensures the
overlap $\abs{c_{0}}^{2}$ is large. This eigenstate-superposition behavior is
precisely what makes QPE useful downstream: in \emph{Shor's algorithm} the
input is an easily prepared superposition over eigenstates of the modular
multiplication operator, and QPE returns a random eigenphase from which the
order is extracted by continued fractions; in \emph{HHL} the input state
$\ket{b}$ is a superposition of eigenvectors of the matrix $A$, and QPE writes
each eigenvalue into an ancilla register so that a controlled rotation can
implement $\lambda\mapsto 1/\lambda$.
\end{remark}

\section{Problems}

\begin{exercise}
Verify the real Hadamard-test identity for a diagonal single-qubit unitary
$U=\operatorname{diag}(1,e^{i2\pi\varphi})$ acting on the eigenstate
$\ket{\psi}=\ket{1}$. Compute $p(0)$ and $p(1)$ directly and confirm they match
Theorem~\ref{l6:thm:hadamard}.
\end{exercise}

\begin{exercise}
Show that the imaginary Hadamard test, obtained by inserting $S^{\dagger}$ on
the ancilla, gives $p(0)=\tfrac12\big(1+\mathrm{Im}\braket{\psi|U|\psi}\big)$.
Explain why the real test alone cannot distinguish $\varphi$ from $1-\varphi$.
\end{exercise}

\begin{exercise}
Let $\varphi=0.011_2=\tfrac38$ and $t=3$. Write down the ancilla state after the
controlled powers in standard QPE, identify it as $\QFT\ket{k}$ for the correct
$k$, and conclude the measurement outcome.
\end{exercise}

\begin{exercise}
For $\varphi=\tfrac13$ (not dyadic) and $t=3$, compute the probabilities
$\abs{\gamma_{k}}^{2}$ for $k=2$ and $k=3$ (the two nearest levels
$\tfrac28=0.25$ and $\tfrac38=0.375$) using
$\gamma_{k}=\tfrac1T\dfrac{1-e^{i2\pi T(\varphi-k/T)}}{1-e^{i2\pi(\varphi-k/T)}}$
with $T=8$, and comment on which outcome is most likely.
\end{exercise}

\begin{exercise}
Using Theorem~\ref{l6:thm:precision}, determine how many ancilla qubits $t$ are
needed to estimate $\varphi$ to $n=6$ bits of precision with success
probability at least $0.9$ (so $\epsilon=0.1$).
\end{exercise}

\begin{exercise}
Explain why Kitaev's bit-by-bit method achieves query complexity
$O(1/\epsilon)$ while the single-qubit estimator needs $O(1/\epsilon^{2})$
samples. In your answer identify precisely which resource (circuit depth versus
number of repetitions) each method spends.
\end{exercise}

\begin{exercise}
Suppose the input to QPE is $\ket{\Psi}=c_{0}\ket{\psi_{0}}+c_{1}\ket{\psi_{1}}$
with $\abs{c_{0}}^{2}=0.8$ and eigenphases $\varphi_{0},\varphi_{1}$ each
exactly representable with the chosen $t$. Describe the joint distribution of
the measured phase and the post-measurement state of the system register.
\end{exercise}

\begin{solution}
\textbf{Exercise 1.} With $\ket{\psi}=\ket{1}$ we have $U\ket{1}
=e^{i2\pi\varphi}\ket{1}$, so $\braket{\psi|U|\psi}=e^{i2\pi\varphi}$ and
$\mathrm{Re}\braket{\psi|U|\psi}=\cos2\pi\varphi$. Following the circuit,
the $\ket{0}$-ancilla branch is $\tfrac12(\ket{1}+e^{i2\pi\varphi}\ket{1})$,
so $p(0)=\tfrac14\abs{1+e^{i2\pi\varphi}}^{2}=\tfrac12(1+\cos2\pi\varphi)$ and
$p(1)=\tfrac12(1-\cos2\pi\varphi)$, matching Theorem~\ref{l6:thm:hadamard}.

\textbf{Exercise 2.} The $S^{\dagger}$ gate multiplies the $\ket{1}$ branch by
$-i$, so the $\ket{0}$-ancilla branch becomes
$\tfrac12(\ket{\psi}-i\,U\ket{\psi})$ and $p(0)=\tfrac14\big(2
+i\braket{\psi|U|\psi}-i\overline{\braket{\psi|U|\psi}}\big)
=\tfrac12(1+\mathrm{Im}\braket{\psi|U|\psi})$. The real test returns only
$\cos2\pi\varphi$, an even function of $\varphi$, so $\varphi$ and $1-\varphi$
give the same $p(0)$; the imaginary test returns $\sin2\pi\varphi$, whose sign
breaks the ambiguity.

\textbf{Exercise 3.} After the controlled powers the ancilla holds
$\tfrac{1}{\sqrt8}\sum_{j=0}^{7}e^{i2\pi(3/8)j}\ket{j}=\QFT\ket{3}$ since
$\varphi=3/8=k/2^{t}$ with $k=3$. The inverse QFT returns $\ket{3}=\ket{011}$
and the measurement reads $011$, i.e.\ $\tilde\varphi=3/8$ with certainty.

\textbf{Exercise 4.} With $T=8$, $\varphi=\tfrac13$. Using
$\abs{\gamma_{k}}^{2}=\tfrac{1}{T^{2}}\frac{\sin^{2}(\pi T(\varphi-k/T))}
{\sin^{2}(\pi(\varphi-k/T))}$: for $k=3$,
$\varphi-k/T=\tfrac13-\tfrac38=-\tfrac{1}{24}$, so
$\abs{\gamma_3}^{2}=\tfrac{1}{64}\frac{\sin^2(\pi/3)}{\sin^2(\pi/24)}
\approx\tfrac{1}{64}\cdot\frac{0.75}{0.01704}\approx0.687$.
For $k=2$, $\varphi-k/T=\tfrac13-\tfrac14=\tfrac{1}{12}$, so
$\abs{\gamma_2}^{2}=\tfrac{1}{64}\frac{\sin^2(2\pi/3)}{\sin^2(\pi/12)}
\approx\tfrac{1}{64}\cdot\frac{0.75}{0.06699}\approx0.175$.
The nearest level $k=3$ ($=0.375$) is most likely, as expected since it is the
closest $t$-bit approximation to $0.333\ldots$

\textbf{Exercise 5.} Take $\epsilon=0.1$, $n=6$. Then
\[
  t=n+\big\lceil\log_2\big(2+\tfrac{1}{2\cdot0.1}\big)\big\rceil
   =6+\lceil\log_2(2+5)\rceil=6+\lceil\log_2 7\rceil=6+3=9
\]
ancilla qubits.

\textbf{Exercise 6.} Kitaev's method spends \emph{circuit depth}: applying
$U^{2^{j}}$ (which costs $2^{j}$ applications of $U$) amplifies the phase so that
each bit is read with high confidence, and the total query cost sums to
$O\big(\sum_j 2^j\big)=O(1/\epsilon)$. The single-qubit estimator uses only $U$
and learns $\varphi$ through a probability that must be resolved to
$O(\epsilon)$; by Bernoulli variance this needs $O(1/\epsilon^2)$
\emph{repetitions}. Depth buys a quadratic saving over samples.

\textbf{Exercise 7.} QPE acts on each eigencomponent independently. With
probability $\abs{c_0}^2=0.8$ the measurement returns $\varphi_0$ exactly (it is
$t$-bit representable) and the system collapses to $\ket{\psi_0}$; with
probability $\abs{c_1}^2=0.2$ it returns $\varphi_1$ and the system collapses to
$\ket{\psi_1}$. The measured phase and the surviving eigenstate are perfectly
correlated: outcome $\varphi_k$ always accompanies $\ket{\psi_k}$.
\end{solution}


\end{document}
